% concatenated from arXiv_CinqueOrsingher_IntegralCompoundProcesses.tex
% INFO DOCUMENTO:


% - - - - -  - - - - -  - - - - -  - - - - -  - - - - -  - - - - -  - - - - -  - - - - - 
% - - - - -  - - - - - INIZIO PREAMBOLO - - - - -  - - - - - 
% - - - - -  - - - - -  - - - - -  - - - - -  - - - - -  - - - - -  - - - - -  - - - - - 

\documentclass[12pt, twoside, english, headsepline]{article}

% - - - - - FORMATTAZIONE - - - - - 
% Link utile: https://www.overleaf.com/learn/latex/page_size_and_margins

\usepackage[margin=30mm]{geometry}
%
% ----- formattazione HEADER AND FOOTER (ricorda `oneside' o `twoside' nella `documentclass')
%\usepackage{scrlayer-scrpage}  % header and footer for KOMA-Script
%\clearscrheadfoot                 % deletes header/footer
%\pagestyle{scrheadings}           % use following definitions for header/footer
%% definitions/configuration for the header
%% r=right, c=center, l=left ; e=even, o=odd
%\cehead[]{F. Cinque and E. Orsingher}        % equal page, right position (inner) 
%\cohead[]{Point processes of the Poisson-Skellam family}        % odd   page, left  position (inner) 
%% definitions/configuration for the footer
%\cefoot[\pagemark]{\pagemark}  % even page, center position
%\cofoot[\pagemark]{\pagemark}     % odd   page, center position
%
%% ----- formattazione TITOLI di sezioni e sottosezioni 
%\usepackage{sectsty}
%\sectionfont{\fontsize{14}{13}\selectfont}
%\subsectionfont{\fontsize{12}{13}\selectfont}

%% - - - - - PACKAGES - - - - - 
%
\usepackage{tikz,graphicx}
\usetikzlibrary{arrows}
\usetikzlibrary{patterns}
\newcommand{\midarrow}{\tikz \draw[-triangle 90] (0,0) -- +(.1,0);}
\tikzset{
  % style to add an arrow in the middle of a path
  mid arrow/.style={postaction={decorate,decoration={
        markings,
        mark=at position .5 with {\arrow[#1]{stealth}}
      }}},
}
\usetikzlibrary{decorations.markings} % per colorare le regioni di un grafico
\tikzset
 {every pin/.style = {pin edge = {<-}}, 
  > = stealth, 
  flow/.style = 
   {decoration = {markings, mark=at position #1 with {\arrow{>}}},
    postaction = {decorate}
   },
  flow/.default = 0.5,   
  main/.style = {line width=1pt}
 }

\usepackage{amsmath,amsthm,amssymb,amsfonts, mathtools, amsbsy, dsfont}
\usepackage[utf8]{inputenx}
%\usepackage[nottoc]{tocbibind}
%\usepackage{cite}
\usepackage{xcolor}
\usepackage{microtype}
\usepackage{hyperref}
\hypersetup{linkcolor=blue}

%\usepackage[pagewise]{lineno}\linenumbers


% - - - - - NEWCOMMAND & COMANDI MATEMATICI - - - - - 

\newcommand*\dif{\mathop{}\!\mathrm{d}}  % simbolo differenziale
\newcommand{\conv}{\mathop{\scalebox{2.8}{\raisebox{-0.2ex}{$\ast$}}}} % convoluzione grande (per fare convoluzioni multipler)
\newcommand{\sconv}{\mathop{\scalebox{2.0}{\raisebox{-0.2ex}{$\ast$}}}} % convoluzione grande (per fare convoluzioni multipler)
\newcommand\sbullet[1][.5]{\mathbin{\vcenter{\hbox{\scalebox{#1}{$\bullet$}}}}} % punto 
\newcommand{\floor}[1]{\left\lfloor #1 \right\rfloor}  % parte intera
\newcommand{\ceil}[1]{\left\lceil #1 \right\rceil}  % parte intera superiore
\newcommand*{\MR}[1]{\href{https://mathscinet.ams.org/mathscinet-getitem?mr=#1}{MR-#1}}  % link MR
\newcommand*{\doi}[1]{\href{http://dx.doi.org/#1}{DOI: #1}}  % link DOI

\allowdisplaybreaks % fa spezzare le formule ``equation'' tra le pagine

% ----- TEOREMI ECC.
\newtheorem{theorem}{Theorem}[section] %[theorem]
\newtheorem{corollary} {Corollary}[section] %[theorem]
\newtheorem{proposition}{Proposition}[section] %[theorem]
\newtheorem{lemma}[theorem]{Lemma}
\theoremstyle{definition} % plain
\newtheorem{definition}{Definition}[section] 
\newtheorem{example}{Example}[section]
%\newtheorem{exercise}{Exercise}[theorem]
\newtheorem{remark}{Remark}[section] %[theorem]
\numberwithin{equation}{section} % equazioni con numero interno alla sezione

% - - - - - KEYWORDS & MSC - - - - - 

% Keywords
\providecommand{\keywords}[1]
{
  \small	
  \textbf{\textit{Keywords: }} #1
}
% MSC
\providecommand{\MSC}[1]
{
  \small	
  \textit{2020 MSC: } #1   
}

% - - - - - TITOLO & AUTORE - - - - - 

\title{Integrals of some compound processes}%Compound processes and their integrals
% RUNNING TITLE: 
\author{Fabrizio Cinque$^1$ and Enzo Orsingher$^2$\\
        \small Department of Statistical Sciences, Sapienza University of Rome, Italy \\
        \small $^1$cinque.fabrizio@gmail.com $^2$enzo.orsingher@uniroma1.it
}

% - - - - - DATA - - - - - 
% inizio 01/04/2021
%\date{March 23, 2025} % Comment this line to show today's date

% - - - - -  - - - - -  - - - - -  - - - - -  - - - - -  - - - - -  - - - - -  - - - - - 
% - - - - -  - - - - - FINE PREAMBOLO - - - - -  - - - - - 
% - - - - -  - - - - -  - - - - -  - - - - -  - - - - -  - - - - -  - - - - -  - - - - - 

% - - - - -  - - - - -  - - - - -  - - - - -  - - - - -  - - - - -  - - - - -  - - - - - 
% - - - - -  - - - - - INIZIO DOCUMENTO - - - - -  - - - - - 
% - - - - -  - - - - -  - - - - -  - - - - -  - - - - -  - - - - -  - - - - -  - - - - - 

\begin{document}

\maketitle

\begin{abstract}
%% Text of abstract
By means of the theory of compound (non-homogenoeus) Poisson processes we define a general framework which includes most of the generalizations of the Poisson processes, both in the Skellam sense and in the space-fractional sense. We prove that Bernstein subordination is the only time-changing leading to a compound Poisson process. We also consider the case of subordination with inverse Bernstein subordinators. Then, we focus on the integrals of compound processes. We give an explicit representation of the fractional integral of compound Poisson process, showing that for fixed $t$ it is distributed as a compound Poisson random variable. We then extend this result to more general integral forms. Furthermore, we obtain some limit results, explicit forms of the iterated integrals and their governing equation. Finally, we study the integral of compound renewal processes in the Fourier-Laplace domain.
\end{abstract} \hspace{10pt}

\keywords{Non-homogeneous Poisson processes; Compound renewal processes; Integral of stochastic processes; Bernstein functions; Convolution-type derivatives}

\MSC{Primary 60G51, 60G55; Secondary 60G22}

%33C10 Bessel and Airy functions, cylinder functions
% 34A05 Explicit solutions, first integrals of ordinary differential equations
%35K25 Higher-order parabolic equations
% 60G55 Point processes
% 60G22 fractional processes, including Brownian motion
% 60G51 processes with independent increments< Levy processes

% ------  CORPO  ---------------------------------------------------------------------------------------

\section{Introduction}

The Poisson process has been generalized in many directions in the recent decades. One first step was obtained by replacing the time-derivative with the time-fractional derivative in the governing equation \cite{L2003}, which first introduced a fractional version of a point process. Another (equivalent) generalization was the replacement of the intertimes between successive events with Mittag-Leffler distributions \cite{MGS2004, MNV2011}. A completely different variation was obtained by acting on the shift space operator in the governing equation \cite{OP2012}, leading to the so-called space-fractional Poisson process and to its generalizations by means of Bernstein subordinators \cite{OT2015}. Another approach was by introducing the generalized counting process where multiple arrivals are permitted \cite{P1984} and by then applying a time-fractional operator to the governing equation \cite{DcMM2016}. Some authors also studied convolutional-type generalization \cite{KM2021}. Parallel to these extensions of the Poisson process,in the literature appeared generalizations of the Skellam process (introduced in \cite{S1946}), obtained by means of the difference of various types of the above Poisson-type processes \cite{BS2024, CO202, KK2024, KV2019, MV2019}.

The first part of this work aims to describe these processes in a more general framework based on the theory of compound Poisson processes, thus assuming jumps of arbitrary real-valued amplitude. Some particular cases of this approach already appeared in the literature, see for instance \cite{BM2012, BM2014, DcMZ2015, KK2023, S2013}. In particular, we study processes (say $M$) with independent increments and having the following infinitesimal behaviour, for $t\ge0, x\in\mathbb{R}$,
\begin{align*}\label{compoundPoissonInfinitesimaleIntroduzione}
	P\{M[t,t+\dif t) \in \dif x\} = \begin{cases}
		\begin{array}{l l}
			\lambda(t)\dif t F(\dif x) + o(\dif t), & x\in\mathbb{R}\setminus\{0\},\\
			1-\lambda(t)\dif t + \lambda(t)\dif t F(\dif 0) + o(\dif t), & x = 0,
		\end{array}
	\end{cases} 
\end{align*}
which prove to be compound (non-homogeneous) Poisson processes. 
%These are among the most popular processes having good theoretical properties and a wide range of possible applications in several different fields.

In the second part of the work we focus on the integral of compound processes (say $M$) of the form $\int_0^t M(s) f(s) \dif s$, $t\ge0$, where $f$ is a suitable real function; in particular we deal with the Poisson and the renewal cases (we refer to \cite{B2022} for compound renewal processes and their asymptotics). For instance, with $t\ge0$ and $f(s) = (t-s)^\alpha,\ \alpha\ge0$, we are able to represent the fractional integral of the compound homogeneous Poisson process $M(t) = \sum_{k=1}^{N(t)}X_k, t\ge0,$ as a compound Poisson random variable,
 $$\int_0^t M(s)(t-s)^\alpha \dif s \stackrel{d}{=} \frac{t^{\alpha+1}}{\alpha+1} \sum_{k=1}^{N(t)} X_k U_k^{\alpha+1},\ \ \ t\ge0,\ \alpha\ge0,$$
 where $U_k\sim Unif(0,1)$ are i.i.d. r.v.s independent from the other random elements.
 
  In the literature, the integration of point processes has been investigated over the years because of its wide applications. Indeed, these integrals are meaningful for inventory and queuing systems and, in general, they can represent the total "cost" of the "elements" released in a time interval by the "subjects" appearing during that interval. For instance, in an enviromental traffic analysis it may represent the total amount of emissions by the cars arriving and stopping at a traffic light; in a biological study, they can represent the food consumption of bacterias or animals arrived in an area and randomly reproducing in a precise time interval. Early works about these processes (and applications) are \cite{GMn1971, Mn1970, N1966, P1966, P1968}, about birth-death process, and \cite{W1971}, about compound Poisson processes. We refer to \cite{N1966} for a discussion on the applications and the study of the moments of integrated Poisson processes. Recently, different authors devoted their research to this type of integral processes, focusing on some path-wise properties \cite{PC2026, PS2002}, on the renewal processes \cite {SVw2007}, on the fractional integrals \cite{CO202, OP2013}, on the Skellam process \cite{X2018}, introducing random times \cite{VK2024}. 
\\



The paper is organized as follows. In the second section we recall some properties of the compound processes and the equality in law in the space of the trajectories of the stochastic processes. In Section 3 we focus on the non-homogeneous Poisson case, showing the connection to the infinitesimal definition of certain point processes. In the homogeneous case we show that this class of processes is closed with respect to the composition with Bernstein subordinators (which produces the so-called Bernstein fractional counting processes \cite{OT2015}). We also briefly study the behaviour of the time-changing with the inverse of a Bernstein subordinator which modifies the compound Poisson process into a compound renewal and we provide the related governing equation.

After some basics on the integral of a stochastic processes, the last section 	of the paper focuses on the integral of compound processes, both Poisson and renewal type. We begin with the analysis of the fractional integral of a compound Poisson process and then extend the result to more general integral forms, proving that for fixed $t$, the integral is distributed as a compound Poisson random variable. We also derive some limit results, the explicit form of the iterated integrals and the equation governing their joint distribution. Then, we include the derivation of the explicit probability law for the integral of Skellam-type processes. At last, we provide some results on the integral of compound renewal processes in the Fourier-Laplace domain.




\section{Preliminaries in compound processes}

Let $(\Omega, \mathcal{F}, \{\mathcal{F}\}_t, P)$ be a filtered probaility space where all the random elements will be defined and  $\mathcal{D}[0,\infty)$ the space of the real c\`{a}dl\`{a}g functions endowed with the Skorokhod topology.

We begin by recalling some known results which will be useful in the rest of the paper.
\begin{lemma}\label{lemmaUguaglianzaProcessiSpazioFunzioniCadlag}
	Let $X, Y$ two stochastic processes on $\mathcal{D}[0,\infty)$. Assume that $X(t)\stackrel{d}{=}Y(t), \ t\ge0$. 
	\begin{itemize}
		\item[($i$)] If both $X$ and $Y$ have independent and stationary increments, then $X\stackrel{d}{=}Y$ on $\mathcal{D}[0,\infty)$.
		\item[($ii$)] If both $X$ and $Y$ have independent increments and $X(t)-X(s)\stackrel{d}{=}Y(t)-Y(s), \ t>s\ge0$, then $X\stackrel{d}{=}Y$ on $\mathcal{D}[0,\infty)$.
	\end{itemize} 
\end{lemma}

We point out that the equality in $\mathcal{D}[0,\infty)$ is important to obtain a pointwise equality for the integral process.

The reader can find different definitions of compound processes (both Poisson's and non), here we use the following one.

\begin{definition}
	Let $N$ be a counting process and $X$ a real r.v. with distribution $F$. We define the compound process of $N$ with $X$ as $M=\big\{M(t) = \sum_{k=1}^{N(t)} X_k\big\}_{t\ge0}$, with i.i.d. $X_k \stackrel{d}{=}X$ and independent of $N$. We write $M\sim Compound(N,X)$ or equivalently $M\sim Compound(N,F)$.	
\end{definition}

If $N$ is a renewal process with waiting times having distribution $G$, then we call $M$ compound renewal process, $M\sim CR(G, X)$ or $M\sim CR(G,F)$. 
If $N$ is a (non-homogeneous) Poisson process with continuous and integrable rate function $\lambda$, denoted by $N\sim PP(\lambda)$, then we call $M$ (non-homogeneous) compound Poisson process, $M\sim CP(\lambda, X)$ or $M\sim CP(\lambda, F)$, having characteristic function
\begin{equation}\label{trasformataFourierProcessoCompound}
	\mathbb{E}e^{i\gamma \big( M(t)-M(s)\big)} = e^{-\int_s^t \lambda(u)\dif u\big(1- \mathbb{E}e^{i\gamma X}\big) }, \ \ \ t> s \ge0,\ \gamma \in \mathbb{R}.
\end{equation}

Hereafter, $F$ is denoting a probability distribution function on $\mathbb{R}$ and, discussing about non-homogenous Poisson processes, we always consider an arbitrary continuous rate function $\lambda$ such that $\int_0^t\lambda(s)\dif s <\infty, \ t\ge0$.


\begin{lemma}\label{lemmaCompostoIncrementiIndipendenti}
	Let $M\sim Compound(N,X)$.
	\begin{itemize}
		\item[($i$)] If $N$ has independent increments, then $M$ has independent increments.
		\item[($ii$)] If $N$ has stationary increments, then $M$ has stationary increments.
	\end{itemize} 
\end{lemma}

\begin{proof}
	Since $X_k$ are i.i.d. and independent of $N$, then $X_{N(t)+k} \stackrel{d}{=} X,\  k\in\mathbb{N},t\ge0$. Basic calculation proves the statement.
\end{proof}

If $X\in\mathbb{N}$ a.s., $M$ is also a counting process and we have the following result.
\begin{proposition}\label{proposizioneCompostoDiCompostoProcessoConteggio}
	Let $N$ be a counting process, $X$ a natural r.v. and $M\sim Compound(N,X)$. 
	%Let $N$ be a counting process, $X$ a natural r.v. and $Y$ a real r.v.. Let $t\ge0$ and define $M(t) = \sum_{k=1}^{N(t)} X_k$, with i.i.d. $X_k \stackrel{d}{=}X$ and independent of $N$. Then, for i.i.d $Y_{k,h},Y_k\stackrel{d}{=}Y,\ k,h\ge1$, independent of $N$ and $X_1,X_2,\dots$,
	 Let $Y$ be a real r.v. and, for $ k=1,\dots,X$, $Y_k$ be an  independent copy of $Y$ (independent of $N$ and $X$). If $Z\sim Compound(M, Y)$ and $W\sim Compound\big(N, \sum_{k=1}^X Y_k\big)$, then $Z(t)-Z(s) \stackrel{d}{=}W(t)-W(s),\ t>s\ge0$.
\end{proposition}

More explicitly, the statement of Proposition \ref{proposizioneCompostoDiCompostoProcessoConteggio} can be expressed as, $t>s\ge0$,
\begin{equation}\label{uguaglianzaDistribuzioneCompostoDiCompostoProcessoConteggio}
	Z(t)-Z(s) = \sum_{k=M(s)+1}^{M(t)} Y_k \stackrel{d}{=} \sum_{k=N(s)+1}^{N(t)} \sum_{h=1}^{X_k} Y_{k,h} = W(t)-W(s),
\end{equation}
where $Y_{k,h}\stackrel{d}{=}Y,\ k,h\ge1$, are independent of $N$ and $X_1,X_2,\dots$.

\begin{proof}
	Thanks to the independence of the elements we have that $Z(t)-Z(s) \stackrel{d}{=} \sum_{k=1}^{M(t)-M(s)} Y_k$ and, for $\gamma\in\mathbb{R}$,
	\begin{align*}
		\mathbb{E}e^{i\gamma \sum_{k=1}^{M(t)-M(s)} Y_k} &= 	\mathbb{E} \Bigl(	\mathbb{E}e^{i\gamma Y} \Big)^{M(t)-M(s)} \\
		& =\mathbb{E} \Bigl( \mathbb{E}e^{i\gamma Y} \Big)^{\sum_{k=1}^{N(t)-N(s)} X_k} \\
		& = 	\mathbb{E} \Bigg( 	\mathbb{E}\Big( 	\mathbb{E}e^{i\gamma Y}\Big)^X \Bigg)^{N(t)-N(s)}.
	\end{align*}
	The proof completes by observing that for $k\ge1, \ \mathbb{E}e^{i\gamma \sum_{h=1}^{X_k} Y_{k,h} } =\mathbb{E}\Big(\mathbb{E}e^{i\gamma Y}\Big)^X$.
\end{proof}


\begin{remark}\label{remarkUguaglianzaSuSpazioFunzioniContinueCompoundDiCompound}
	By keeping in mind the hypotheses of Proposition \ref{proposizioneCompostoDiCompostoProcessoConteggio}, this result, together with Lemma \ref{lemmaUguaglianzaProcessiSpazioFunzioniCadlag} and Lemma \ref{lemmaCompostoIncrementiIndipendenti}, imply that if $N$ has independent increments, then $Z$ has independent increments and $Z\stackrel{d}{=}W$ on $\mathcal{D}[0,\infty)$. Furthermore, if $N$ has also stationary increments then $Z$ ia a L\'evy process (for point ($ii$) of Lemma \ref{lemmaIncrementiIntegraleProcesso}). \hfill $\diamond$
\end{remark}


\section{Compound Poisson process}


\begin{theorem}
	$M\sim CP(\lambda, F)$ if and only if $M$ is a stochastic process having independent increments, $M(0)=0$ a.s. and, for $t\ge0, x\in\mathbb{R}$,
	\begin{align}\label{compoundPoissonInfinitesimale}
		P\{M[t,t+\dif t) \in \dif x\} = \begin{cases}
			\begin{array}{l l}
				\lambda(t)\dif t F(\dif x) + o(\dif t), & x\in\mathbb{R}\setminus\{0\},\\
				1-\lambda(t)\dif t + \lambda(t)\dif t F(\dif 0) + o(\dif t), & x = 0.
			\end{array}
		\end{cases} 
	\end{align}
\end{theorem}

Note that $F(\dif x) = F(x)-F(x^-) = p(t, \dif x)$, i.e. if $F$ is absolutely continuous, $F(\dif x) = f(x)\dif x$, otherwise it represents the mass probability at point $x$. The same holds for $p(t,\dif x) =P\{M(t) \in \dif x\},\ t\ge0,x\in\mathbb{R}$.

\begin{proof}
	Sufficiency. The non-homogeneous Poisson process $N$ with rate function $\lambda$ has independent increments and $M\sim CP(\lambda, F)$ as well since ($i$) of Lemma \ref{lemmaCompostoIncrementiIndipendenti}. Obviously, $M(0) = 0$. By keeping in mind that $X_1,X_2,\dots \sim F$ are i.i.d. and independent of $N$, we have that
\begin{align}
	P\{M[t,t+\dif t)\in \dif x\} %&= \sum_{n\ge0} P\Big\{\sum_{k=N(t)}^{N[t,t+\dif t)} X_k\in \dif x, N[t, t+\dif t) = n\Big\} \\
	& = \sum_{n=0,1} P\left\{ \sum_{k=1}^{N[t,t+\dif t)} X_k\in \dif x, N[t,t+\dif t) = n\right\} \nonumber \\
	& = \big(1-\lambda(t)\dif t\big) \delta_0(\dif x) + \lambda(t)\dif t P\{X\in \dif x\} + o(\dif t), \label{dimostrazioneRappresentazioneInfinitesimaleCompoundPoisson}
\end{align}	
with $\delta_0$ being the Dirac measure. Formula \eqref{dimostrazioneRappresentazioneInfinitesimaleCompoundPoisson} coincides with \eqref{compoundPoissonInfinitesimale}.
	
	Necessity. By means of usual arguments we obtain 
% that,	$$p(t + \dif t, \dif x) = \lambda(t) \int_{\mathbb{R}} p(t, \dif x-y)F(\dif y) + p(t, \dif x)\big(1-\lambda(t) \dif t\big)$$
	the governing forward equation
	\begin{align}
		\frac{\dif}{\dif t}p(t, \dif x) &= \lambda(t) \Bigg( \int_{\mathbb{R}} p(t, \dif x-y)F(\dif y) - p(t,\dif x)\Bigg) \nonumber\\
		& = \lambda(t) \big( F \ast p (t, \dif x) - p(t, \dif x) \big),\ \ \ x\in\mathbb{R}, \label{equazioneCompoundPoisson}
	\end{align}
	subject to $p(0, \dif x) = \delta_0(\dif x),\ x\in\mathbb{R},$.
	\\Now, by means of the $x$-Fourier transform and using the notation $\mathcal{F} f(\gamma) = \int_\mathbb{R} e^{i\gamma x}f(x)\dif x,\ \gamma\in\mathbb{R}$, equation \eqref{equazioneCompoundPoisson} turns into
	\begin{equation}
		\frac{\dif}{\dif t}\mathcal{F}p(t, \gamma) = \lambda(t)\Big(\mathbb{E}e^{i\gamma X} -1\Big)\mathcal{F}p(t,\gamma),\ \ \ \mathcal{F}p(0,\gamma) = 1,\ \ \ \gamma \in\mathbb{R},
	\end{equation}
	and \eqref{trasformataFourierProcessoCompound} with $s=0$ solves it. The proof concludes by using point ($ii$) of Lemma \ref{lemmaUguaglianzaProcessiSpazioFunzioniCadlag} and the independence of the increments, which implies that $M(t)-M(s)$ has characteristic function \eqref{trasformataFourierProcessoCompound}.
\end{proof}

\begin{remark}
	We point out that in the proof above we have showed that the probability law of $M\sim CP(\lambda, F)$ satisfies equation \eqref{equazioneCompoundPoisson}, which equivalently writes
	\begin{equation}\label{equazioneOperatorialeCompoundPoisson}
		\frac{\dif}{\dif t}p(t, \dif x) = -\lambda(t) \big( I- C_F \big) p (t, \dif x),\ \ \ t\ge0,\ x\in\mathbb{R},
	\end{equation}
	where $I$ is the identity operator and $C_F g(t, x) = F(t,x)\ast g = \int_\mathbb{R} g(t, x-y) F(\dif y)$ denotes the convolution operator with respect to the distribution $F$. In the homogeneous case, we are using this formulation to study the effect of applying a Bernstein function to the operator on the right-hand side of \eqref{equazioneOperatorialeCompoundPoisson}, see Theorem \ref{teoremaComposizioneCompoundPoissonConSubordinatore}. \hfill$\diamond$
\end{remark}

\begin{example}[Skellam-type processes]\label{esempioSkellam}
	Let $\mathcal{I}\subset \mathbb{R}\setminus\{0\}$ be a finite set and $M\sim CP(\lambda, F)$ with distribution $F(x) = \sum_{i\in \mathcal{I}} a_i \mathds{1}(x\ge i)$, such that $a_i>0\ \forall\ i,\ \sum_{i\in\mathcal{I}}a_i = 1$. Thus, $F(\dif x)= \sum_{i\in\mathcal{I}} a_i \delta_i(\dif x)$, where $\delta_i$ is the Dirac delta measure centred in $i$. $M$ is a discrete point process which jumps, at Poisson times, with size $i\in \mathcal{I}$ with probability $a_i$. This is a Skellam-type process, see for instance \cite{CO202} and references therein; if $\mathcal{I} = \{-K,\dots,-1,1,\dots,K\}$, with $K\in\mathbb{N}$, $M$ is called a Skellam process of order $K$; if $\mathcal{I}=\{1\}$ then $M$ is a Poisson process. 
	
	Hence, if $M\sim CP(\lambda, F)$ with $F$ defined above and  $N_i\sim PP(a_i\lambda) = CP(a_i\lambda, 1), \ i\in\mathcal{I},$ are independent Poisson processes, then  $M\stackrel{d}{=} \sum_{i \in\mathcal{I} } i N_i.$
	
	From \eqref{equazioneCompoundPoisson} we easily derive the differential equation governing the probability law of $M$. For $x\in\mathbb{Z}$,
	\begin{align}
			\frac{\dif}{\dif t}p(t, x) &= \lambda(t) \Bigg( \int_{\mathbb{R}} p(t, x-y)\sum_{i\in \mathcal{I}} a_i\delta_i(\dif y) - p(t, x)\Bigg)\nonumber\\
			& =\lambda(t)\sum_{i\in\mathcal{I}}a_i \big(p(t, x-i) -p(t,x)\big),\ \ \ t\ge0,
	\end{align}
	which is a particular case of formula (2.4) of \cite{CO202} (with $\lambda_i = a_i\lambda$). \hfill$\diamond$
\end{example}




\subsection{Homogeneous case}

We recall that $f:[0,\infty)\longrightarrow [0,\infty)$ is a Bernstein function if $f\in C^\infty, \ (-1)^n \dif^n f/\dif x^n \le 0, \ n\ge1$, and it can be expressed as
\begin{equation}\label{definizioneFunzioneBernstein}
	f(x) = a+bx+ \int_0^\infty\Bigl( 1-e^{-xw}\Bigr) \bar \nu_f(\dif w), \ \ \ x\ge0,
\end{equation}
where $a,b\ge0$ and $\bar \nu_f$ is a L\'{e}vy measure, i.e. such that $\int_0^\infty (s\wedge1)\bar \nu_f(\dif s)<\infty$. We refer to \cite{SSV2012} for a detailed discussion.
Bernstein functions are related to non-decreasing L\'{e}vy processes, also known as subordinators. Indeed, for each Bernstein function $f$ there exists a subordinator $\mathcal{H}_f$ such that $f$ is the L\'{e}vy symbol of $\mathcal{H}_f$, i.e. $\mathbb{E}e^{-\mu\mathcal{H}_f(t)} = e^{-tf(\mu)}, \ \mu, t\ge0$ % Hereafter we assume $a=b=0$
(and, as well-known, $\mathcal{H}_f$ has independent and stationary increments).


\begin{lemma}\label{lemmaEquazioneOperatorialeBernstein}
	Let $\lambda>0$, $f$ be a Bernstein function as in \eqref{definizioneFunzioneBernstein} with $ \bar \nu_f\not=0$ and
	\begin{equation}\label{pesiBernstein}
		a_k = a_{\lambda,f,k} = \frac{1}{f(\lambda)} \frac{\lambda^k}{k!}\int_0^\infty e^{-\lambda w} w^k \bar \nu_f(\dif w),\ \ \ k\ge1.
	\end{equation}
	\begin{itemize}
		\item[($i$)] $\displaystyle \sum_{k\ge1} a_k = 1-(a+b\lambda)/f(\lambda)$.
		\item[($ii$)] $ \displaystyle f\big(\lambda(1-x)\big) = f(\lambda) \Bigg(1- \sum_{k\ge 1}a_k x^k\Bigg) - b\lambda x ,\ \ \ x\le1. $
		\item[($iii$)] Let  $a'_1 = a_1 + b\lambda/f(\lambda)$ and $a'_k = a_k,\ k\ge2$. Let $F$ be a distribution function and $p$ be a suitable function. $f$ has $a=0$ if and only if
		\begin{equation}\label{relazioneOperatorialeFunzioneBernstein}
			f\Big(\lambda\big(I-C_F\big)\Big) p= f(\lambda)\big(I-C_{ F_f}\big) p
		\end{equation}
		with $F_f$ being the distribution
		\begin{equation}\label{distribuzioneComposizioneBernstein}
		%	F_f(x) = \frac{1}{f(\lambda)}\sum_{k\ge1} \frac{\lambda^k}{k!}\int_0^\infty e^{-\lambda w} w^k \bar \nu_f_f(\dif w) \, F^{(k)}(x), \ \ \ x\in \mathbb{R}
		F_f = \sum_{k\ge1} a'_k F^{(k)} = \sum_{k\ge1} a'_k\, \conv_{h=1}^k F, \ \ \ x\in \mathbb{R}.
		\end{equation}
	\end{itemize} 
\end{lemma}

%Hereafter we denote the weights of the combination in \eqref{distribuzioneComposizioneBernstein} by $a_k,\ k\ge1$, i.e. $ F_f = \sum_{k \ge1} a_k F^{(k)}$.

\begin{proof}
	($i$) For any $a,b\ge0$, we have
	\begin{align}
		\sum_{k\ge1} a_k % \frac{1}{f(\lambda)}\sum_{k\ge1} \frac{\lambda^k}{k!}\int_0^\infty e^{-\lambda w} w^k \bar \nu_f(\dif w) 
		&= \frac{1}{f(\lambda)}\int_0^\infty e^{-\lambda w}\sum_{k\ge1} \frac{(\lambda w)^k}{k!} \bar \nu_f(\dif w)\nonumber \\
		& =\frac{1}{f(\lambda)} \int_0^\infty  e^{-\lambda w} (e^{\lambda w}-1)\bar \nu_f(\dif w) \nonumber\\
		& = \frac{f(\lambda)-a- b\lambda}{f(\lambda)}.\label{calcoliPesiBernstainSommano1}
	\end{align}
	($ii$) For any $a,b\ge0$ and $x\le 1$, we have
	\begin{align}
		f\big(\lambda(1-x)\big) & = a+b\lambda(1-x) +\int_0^\infty \Big(1-e^{-\lambda w} \sum_{k\ge0}\frac{(\lambda x w)^k}{k!} \Big)\bar \nu_f(\dif w)\nonumber\\
		& =  a + b\lambda(1-x) + \int_0^\infty (1-e^{-\lambda w})\bar \nu_f (\dif w) - \sum_{k\ge1}\frac{(\lambda x )^k}{k!}  \int_0^\infty e^{-\lambda w} w^k \bar \nu_f(\dif w)\nonumber\\
		& = f(\lambda)\Big(1- \frac{1}{f(\lambda)}\sum_{k\ge1}\frac{(\lambda x )^k}{k!}  \int_0^\infty e^{-\lambda w} w^k \bar \nu_f(\dif w)\Big) -b\lambda x.
	\end{align}
	($iii$) $F_f$ is a distribution since it is a convex combination of the distributions $F^{(k)},\ k\ge1$. Indeed, for point ($i$) we have $\sum_{k} a'_k = \sum_{k}a_k+\lambda b/f(\lambda) = 1-a/f(\lambda)$ which is $1$ if and only if $a=0$.
	Now, by suitably adapting the steps of point ($ii$) we obtain
	\begin{align}
		f\Big(\lambda\big(I-C_F\big)\Big) p& =  f(\lambda)\Big(I- \frac{1}{f(\lambda)}\sum_{k\ge1}\frac{\lambda ^k}{k!}  \int_0^\infty e^{-\lambda w} w^k \bar \nu_f(\dif w)C_F^k\Big)p - b\lambda C_Fp \nonumber \\
		& =  f(\lambda)\Bigg(p- \Big(a_1+\frac{\lambda b}{f(\lambda)}\Big)F\ast p - \sum_{k\ge2}a_k \Big(\conv_{h=1}^k F\Big)\ast p \Bigg) \nonumber \\
		& = f(\lambda) \Big(p- \sum_{k\ge1} a'_k F^{(k)}\ast p \Big)\nonumber\\
		& =  f(\lambda) \Big(I- C_{F_f} \Big)p \nonumber
	\end{align}
	which proves \eqref{relazioneOperatorialeFunzioneBernstein}.
\end{proof}

%\begin{theorem}\label{teoremaEquazioneDifferenzialeBernsteinCompound}
%	Let $\lambda>0$, $F$ be a distribution function and $f$ be a Bernstein function with $a=b=0, \bar \nu_f\not=0$. Then, the solution to 
%		\begin{equation}\label{equazioneDifferenzialeOperatorialeFunzioneBernstein}
%		\frac{\dif}{\dif t}p(t, \dif x) = - f\Big(\lambda\big(I-C_F\big)\Big)p(t, \dif x) 
%	\end{equation}
%	is the probability mass function of $M_f\sim CP\big(f(\lambda), F_f\big)$, with $F_f$ in \eqref{distribuzioneComposizioneBernstein}, and $M_f = M\circ\mathcal{H}_f$. Finally, $M_f$ has jumps distributed as
%	\begin{equation}\label{leggeSaltiCompoundBernstein}
%		Y = \sum_{k\ge1} \mathds{1}(B = k) \sum_{h=1}^k X_h,\ \ \text{such that}\ \ P\{B = k\} = a_k,\ k\ge1,
%	\end{equation}
%	with $B$ and $X_1,X_2,\dots\sim F$ independent.
%\end{theorem}
%
%\begin{proof}
%	
%	The jumps \eqref{leggeSaltiCompoundBernstein} follow from the structure of the distribution $F_f$ in \eqref{distribuzioneComposizioneBernstein}.
%\end{proof}


\begin{theorem}\label{teoremaComposizioneCompoundPoissonConSubordinatore}
Let $M\sim CP(\lambda, F)$, $f$ be a Bernstein function (see \eqref{definizioneFunzioneBernstein}) and $\mathcal{H}_f$ be the corresponding Bernstein subordinator.
\begin{itemize}
	\item[($i$)] $a=0$ if and only if $M\circ \mathcal{H}_f\sim CP\big(f(\lambda), F_f\big)$, with %$F_f = F$ if $b\not=0$ and 
	$F_f$ in \eqref{distribuzioneComposizioneBernstein}. The jumps are distributed as $Y = \sum_{h=1}^A X_h$, with $A$ being independent of $X_1,X_2,\dots\sim F$ and $M$ and $P\{A = k\} = a'_k,\ k\ge1,$ with $a'_k$ in ($iii$) of Lemma \ref{lemmaEquazioneOperatorialeBernstein}.
	%\begin{equation}\label{leggeSaltiCompoundBernstein}
	%Y = \sum_{k\ge1} \mathds{1}(B = k) \sum_{h=1}^k X_h,\ \ \text{such that}\ \ P\{B = k\} = a_k,\ k\ge1,
	%\end{equation}
	\item[($ii$)] Let $a=0$. The probability law of $M\circ\mathcal{H}_f$ satistfies, for $x\in\mathbb{R}$,	\begin{equation}\label{equazioneDifferenzialeOperatorialeFunzioneBernstein}
		\frac{\dif}{\dif t}p(t, \dif x) = - f\Big(\lambda\big(I-C_F\big)\Big)p(t, \dif x), \ t\ge0,\ \ \ p(0,\dif x) = \delta_0(\dif x).
	\end{equation}
\end{itemize}
\end{theorem}

Note that statement ($i$) would actually need the hypothesis $f\not=0$ in order to avoid the degenerate case in which $f$ is identically null (which is of no interest and therefore omitted).

\begin{proof}
	($i$) We use the characteristic function of $M\circ \mathcal{H}_f$. Let $\gamma \in \mathbb{R}$ and $t>s\ge0$, 
	\begin{align}
		\mathbb{E} e^{i\gamma \big(M(\mathcal{H}_f(t))-M(\mathcal{H}_f(s))\big)} &= \mathbb{E} \exp\Big(-\lambda \big(\mathcal{H}_f(t)-\mathcal{H}_f(s)\big) \big(1-\mathbb{E}e^{i\gamma X}\big) \Big)\nonumber\\
		& = \exp\Bigg(-(t-s)f\Big(\lambda\big(1-\mathbb{E}e^{i\gamma X}\big) \Big) \Bigg).\nonumber
	\end{align}
	Thus, $M\circ \mathcal{H}_f$ has exponential characteristic function and stationary increments and therefore it also has independent increments. Thus, it sufficies to prove that $M\circ \mathcal{H}_f$ is pointwise compound Poisson. For $t\ge0$,
		\begin{align}
		\mathbb{E} e^{i\gamma M\big(\mathcal{H}_f(t)\big)} &= \exp\Bigg(-t f\Big(\lambda\big(1-\mathbb{E}e^{i\gamma X}\big) \Big) \Bigg) \nonumber\\
		& = \exp \Bigg(-f(\lambda)\Big(1-\sum_{k\ge1} a_k\big(\mathbb{E}e^{i\gamma X}\big)^k \Big) - b\lambda \mathbb{E}e^{i\gamma X} \Bigg) \nonumber\\
		& =  \exp \Bigg(-f(\lambda)\Big(1-\sum_{k\ge1} a'_k \mathbb{E}e^{i\gamma \sum_{h=1}^k X_h}\Big) \Bigg),
	\end{align}
	with $X_1,\dots $ indpendent copies of $X$. In the second step we used ($ii$) of Lemma \ref{lemmaEquazioneOperatorialeBernstein}. Now, since $\sum_{k} a'_k = 1$ if and only if $a=0$ we obtain that $\sum_{k\ge1} a'_k  \mathbb{E}\exp\big(i\gamma \sum_{h=1}^k X_h\big) = \mathbb{E}\exp\big(i\gamma \sum_{h=1}^A X_h\big)$ with $A$ independent r.v. such that $P\{A = k\}=a'_k,\ k\ge1$. Thus, $M\circ\mathcal{H}_f\sim CP\big(f(\lambda), F_f\big)$, where $F_f$ derives from the structure of the jumps, see also \eqref{distribuzioneComposizioneBernstein}.
	\\
	($ii$) From point ($i$) and \eqref{equazioneOperatorialeCompoundPoisson}, $x\in\mathbb{R}$,
	\begin{equation*}
		\frac{\dif}{\dif t}p(t, \dif x) = - f(\lambda)\big(I-C_{F_f}\big)p(t, \dif x), \ t\ge0,\ \ \ p(0,\dif x) = \delta_0(\dif x).
	\end{equation*} 
	The statement follows from point ($iii$) of Lemma \ref{lemmaEquazioneOperatorialeBernstein}.	
\end{proof}



\begin{theorem}\label{teoremaComposizioneCompoundPoissonGenerale}
	Let $M\sim CP(\lambda, F)$ and $T$ be a non-decreasing stochastic process such that $T(0) = 0$ a.s.. $M\circ T\sim CP$ if and only if $T$ is a Bernstein subordinator ($\mathcal{H}_f$) with $f$ having $a=0$. 
\end{theorem}

\begin{proof}
	Necessity. It is point ($i$) of Theorem \ref{teoremaComposizioneCompoundPoissonConSubordinatore}.
	\\
	Sufficiency. Let $\gamma\in\mathbb{R}$,
	\begin{equation}
 		\mathbb{E} e^{i\gamma \big(M(T(t))-M(T(s))\big)} = \mathbb{E} \exp\Big(-\lambda \big(T(t)-T(s)\big) \big(1-\mathbb{E}e^{i\gamma X}\big) \Big).\nonumber
	\end{equation}
	By recalling form \eqref{trasformataFourierProcessoCompound} of the characteristic function of a compound Poisson process, $T$ must have stationary increments and exponential characteristic function. Thus, $T$ must also have independent increments and it can only be a Bernstein subordinator (with $a=0$).
\end{proof}

\begin{example}[Space-fractional Poisson process]
	As an example, we show that if $M$ is a homogeneous Poisson process and $\mathcal{H}_f$ is a Bernstein subordinator with $a=b=0$, Theorem \ref{teoremaComposizioneCompoundPoissonConSubordinatore} easily yields the process defined in equation (1.1) of \cite{OT2015}. Indeed, $M\sim PP(\lambda) = CP(\lambda, 1)$, then $F(\dif x) = \delta_{1}(\dif x)$ and, from \eqref{distribuzioneComposizioneBernstein}, $F_f(\dif x) = \sum_{k\ge1} a'_k \sconv_{h=1}^k F(\dif x) = \sum_{k\ge1} a_k \delta_k(\dif x)$ (note that $a'_k = a_k\ \forall\ k$ since $b=0$, see point ($iii$) of Lemma \ref{lemmaEquazioneOperatorialeBernstein}). Therefore, Theorem \ref{teoremaComposizioneCompoundPoissonConSubordinatore} implies that $N\circ \mathcal{H}_f \sim CP\big(f(\lambda), F_f\big)$, that is a process having infinitesimal behavior (see \eqref{compoundPoissonInfinitesimale}), for $t\ge0$,
	\begin{align*}
		P\{M[t,t+\dif t)  = x\} = \begin{cases}
			\begin{array}{l l}
				f(\lambda)\dif t \,a_x + o(\dif t), & x\in\mathbb{N},\\
				1-f(\lambda)\dif t + o(\dif t), & x = 0,\\
				o(\dif t), & otherwise,
			\end{array}
		\end{cases}
	\end{align*}
	which coincides with (1.1) of \cite{OT2015} by using \eqref{pesiBernstein}. \hfill$\diamond$
\end{example}

We now study the time-changed compound Poisson process with the inverse $L_f$ of the Bernstein subordinator $\mathcal{H}_f$, i.e. $L_f(t)=\inf\{s\ge0\,:\, \mathcal{H}_f(s)\ge t\},\ t\ge0$.
%It is well-known that 
%\begin{equation}
%	\int_0^\infty e^{-\mu t} l_f(t,s)\dif t = -\frac{1}{\mu}  \frac{\partial }{\partial s} \mathbb{E} e^{-\mu \mathcal{H}_f(s)}= \frac{f(\mu)}{\mu} e^{-s f(\mu)},\ \ \ \mu>0.
%\end{equation}

We recall (see \cite{K2011}) that if $N\sim PP(\lambda)$ and $L_f$ is an independent inverse subordinator, then, $N\circ L_f$ is a renewal process with waiting times $Y$ s.t. 
	\begin{equation}\label{intertempoComposizionePoissonInversoSubordinatore}
		P\{Y\le t\} =  1-\mathbb{E}e^{-\lambda L_f(t)}\ \text{ and } \ \mathbb{E}e^{-\mu Y} = \frac{\lambda}{\lambda+f(\mu)},\ \mu, t>0.
	\end{equation}
Thus, if $M\sim CP(\lambda, X)$, then $M\circ L_f\sim CR(Y, X)$ with $Y$ distributed as described in \eqref{intertempoComposizionePoissonInversoSubordinatore}.

%\begin{proof}
%	We begin by showing that the first intertime ($W_1$) of the time-changed process is distributed as the $X$ in the statement. For $t\ge0$,
%	\begin{align}
%		P\{W_1 < t\} = 1-P\{N\circ L_f(t) = 0\} = 1-\mathbb{E}e^{-\lambda L_f(t)},
%	\end{align}
%	and its Laplace transform is
%	\begin{align}
%		\mathbb{E}e^{-\mu W_1} = 
%	\end{align}
%\end{proof}


Now, we show the equation governing the probability law of the compound renewal process $M\circ L_f$, which requires the convolutional-type derivative introduced by different authors, see \cite{K2011, T2015}, 
\begin{equation}\label{definizioneDervitaConvoluzionale}
	\mathcal{D}_t^f u(t) = b\frac{\dif }{\dif t}u(t) + \int_0^t \frac{\partial }{\partial t}u(t-s)\nu_f(s)\dif s,
\end{equation}
where $\nu_f$ is the tail of the L\'{e}vy measure $\bar \nu_f$. We refer to Section 2 of \cite{BS2024} (or Section 4 of \cite{CO2026}) for some information on the operator in \eqref{definizioneDervitaConvoluzionale} and to \cite{K2011, T2015} for further details.
\\

Hereafter we assume the following condition.
\\
\textbf{Condition I.} The L\'{e}vy measure $\bar\nu_f$ associated to the Bernstein function $f$ is such that $\bar\nu(0,\infty) = \infty$ and its tail $\nu_f(s) = a+\bar \nu_f(s,\infty)$ is absolutely continuous.
\\

It was shown in \cite{T2015} that under Condition I, the inverse $L_f$ admits a probability density $l_f(t,x) = P\{L_f(t)\in\dif x\}/\dif x, \ x,t \ge0,$ with $l_f(0,x) = \delta_0(x)$.

In order to prove the next statement we also recall that in \cite{T2015} the author introduced another convolution-type derivative generalizing the Riemann-Liouville fractional derivative, 
\begin{align}
	\mathbb{D}_t^f u(t) =  b\frac{\dif }{\dif t}u(t) + \frac{\dif }{\dif t}\int_0^t u(t-s)\nu_f(s)\dif s,
\end{align}
which relates to the convolution-type derivative \eqref{definizioneDervitaConvoluzionale} by means of (see Proposition 2.7 of \cite{T2015})
\begin{align}\label{relazioneOperatoriConvoluzione}
	\mathcal{D}_t^f u(t) =\mathbb{D}_t^f u(t) - \nu_f(t) u(0).
\end{align}
Furthermore, it is useful to recall (see Theorem 4.1 of \cite{T2015}) that the probability density $l_f$ satisfies the problem, with  $x>0$ if $b = 0$ or $0<x<t/b$ if $b>0$,
\begin{align}\label{legameDerivataConvoluzionaleLiouvilleOperatoreSpaziale}
	\mathbb{D}_t^f l_f(t,x)= -\frac{\partial }{\partial x}l_f(t,x), \ \ \ t\ge0,
\end{align}
subject to $l_f(t,0) = \nu_f(t), \ l_f(t, t/b)=0$ and $l_f(0,x) = \delta(x)$.

\begin{proposition}[Time-fractional compound Poisson process]\label{proposizioneComposizioneInversoSubordinatore}
	Let $M\sim CP(\lambda, F)$, $L_f$ be the inverse Bernstein subordinator and  $p_f$ be the probability function of $M\circ L_f$. Let assume Condition I and that $\mathcal{D}_t^f p_f(t,\dif x) = \int_0^\infty p(s,\dif x) \mathcal{D}_t^f l_f(t,s)\dif s\ \forall t,x$. Then, $p_f$ satisfies
	\begin{equation}
		\mathcal{D}_t^f p_f(t,\dif x) = -\lambda(I-C_F) p_f(t,\dif x),\ \ \ t\ge0,\ x\in \mathbb{R},
	\end{equation}
	subject to $p_f(0,\dif x) = \delta_0(\dif x)$.
\end{proposition}

\begin{proof}
	It is easy to evaluate the initial condition $p_f(0,\dif x) = \int_0^t p(s, \dif x)\delta_0(s) =  p(0,\dif x) = \delta_0(\dif x)$. Then,
	\begin{align*}
		\mathcal{D}_t^f p_f(t, \dif x) &= \int_0^\infty p(s,\dif x) \mathcal{D}_t^f l_f(t,s)\dif s \nonumber\\
		& =  \int_0^\infty p(s,\dif x) \Big( \mathbb{D}_t^f l_f(t,s) -\nu_f(t) l_f(0,s) \Big)\dif s \nonumber\\
		& =  \int_0^\infty p(s,\dif x) \mathbb{D}_t^f l_f(t,s) - p(0,\dif x)\nu_f(t) \nonumber\\
		& =  -\int_0^\infty p(s,\dif x) \frac{\partial}{\partial s}l_f(t,s)\dif s - \delta_0(\dif x) \nu_f(t)\nonumber\\
		& = \int_0^\infty \frac{\partial}{\partial s} p(s,\dif x)  l_f(t,s)\dif s\\
		&= - \lambda \int_0^\infty (I-C_F)p(s,\dif x)  l_f(t,s)\dif s
	\end{align*}
	where in the second-last step we integrated by parts. The proof completes by observing that we can exchange the integral and the operator $(I-C_F)$ since $C_F$ is the convolution with respect to a distribution $F$. 
\end{proof}

%The proof uses the same arguments presented in the proof of Theoerem 4.1 of \cite{CO2026} and therefore it is omitted (also see Theorem 1 of \cite{BS2024}). 

We point out that Proposition \ref{proposizioneComposizioneInversoSubordinatore} is a particular case of Theorem 4.1 of \cite{CO2025}. Indeed, it sufficies to observe that $l_f$ satisfies problem (4.4) of \cite{CO2025} (restricted to $x>0$) suitably adapted to the operator \eqref{definizioneDervitaConvoluzionale} and that this satisfies hypothesis (4.1) of \cite{CO2025}.
%$$\int_0^\infty e^{-\mu t} \mathcal{D}_t^f u(t,x)\dif x = .$$


\section{Integral of compound processes}

Let $T>0$, $X=\{X_n\}_{n\in\mathbb{N}_0}$ be a discrete real stochastic process, $N$ be a counting process wih arrival times $S_k,\ k\ge1$ ($S_0 = 0$ a.s), and measurable $f:[0,\infty)\longrightarrow\mathbb{R}$, $f\in \mathcal{L}^1[0,T]$. In the present section we assume the following condition on $X$ and $N$.
\\

\textbf{Condition II.} If $X$ is a bounded process or $X_n\ge0 (\le0)\ \forall\ n$ a.s. and $f\ge0(\le0)$, then we assume that $N$ can explode in finite time intervals. Otherwise, we assume that $N$ cannot explode in finite intervals.
\\

We are interested in integrals over $[0,t],\ 0\le t\le T,$ which, in light of Condition II, can be written as (see Appendix \ref{appendiceCalcoloEsplicitoIntegrali} for details)
\begin{equation}\label{formulazioneGeneraleIntegraleProcessoTempoContinuo}
\int_0^t X_{N(s)} f(s)\dif s = \sum_{n=0}^{N(t)-1} X_n\int^{S_{n+1}}_{S_n} f(s)\dif s +X_{N(t)}\int_{S_{N(t)}}^t f(s)\dif s.
\end{equation}

If $X$ is a random walk, i.e. $X_n = \sum_{k=1}^n Y_k,\ n\in\mathbb{N}_0,$ with $Y_k$ i.i.d. and $N$ is independent of $Y_k\ \forall\ k$, then $X_{N}$ is a compound process and the integral \eqref{formulazioneGeneraleIntegraleProcessoTempoContinuo} turns into
\begin{equation}\label{formulazioneGeneraleIntegraleProcessoCompound}
 \int_0^t \sum_{k=1}^{N(s)} Y_k f(s) \dif s = \sum_{k=1}^{N(t)} Y_k \int^t_{S_k} f(s) \dif s,\ \ \ 0\le t\le T,
\end{equation}
see Appendix \ref{appendiceCalcoloEsplicitoIntegrali} for some details.

Note that, the characteristic function of the integral \eqref{formulazioneGeneraleIntegraleProcessoCompound} can be written as follows. By taking into account the independence of the $Y$s, we have that, with $\gamma\in\mathbb{R}$, %(with $\mathbb{E}_Y$ denoting the expectation wiht respect to the $Y$s only),
\begin{align}
	\mathbb{E} e^{ i\gamma \int_0^t M(s)  \dif s} &= \sum_{n\ge0} P\{N(t) = n\} \mathbb{E}\Bigg[\prod_{k=1}^n\mathbb{E} \Big[ e^{i\gamma Y_k \int_{S_k}^tf(s)\dif s} \,\Big|\,S_1,\dots,S_n\Big]\,|\,N(t) = n\Bigg]\nonumber\\
	& =  \sum_{n\ge0} P\{N(t) = n\} \mathbb{E}\Bigg[\prod_{k=1}^n H_Y\bigg(\gamma \int_{S_k}^tf(s)\dif s\bigg) \,|\,N(t) = n\Bigg] \label{funzioneCaratteristicaIntegraleGenerale}
\end{align}



If $f\equiv 1$, we recover the following result on the increments of the integral.

\begin{lemma}\label{lemmaIncrementiIntegraleProcesso}
	Let $X$ be a stochastic process on $\mathcal{D}[0,T]$ such that $X(0) =0$ a.s.% and integrable over $[0,T]$
	. Let $I(t) = \int_0^t X(s)\dif s,\ 0\le t\le T$, and $0\le t_1<t_2\le T, \ \gamma \in\mathbb{R}$.
	\begin{itemize}
		\item[($i$)] If $X$ has independent increments,
		\begin{equation}\label{incrementoIntegraleProcessoStocasticoIncrementiIndipendenti}
			\mathbb{E}e^{i\gamma \big(I(t_2)-I(t_1)\big)} = \mathbb{E}e^{i\gamma (t_2-t_1)X(t_1)}\mathbb{E}e^{i\gamma \int_0^{t_2-t_1} \big(X(s+t_1)-X(t_1)\big)\dif s}.
		\end{equation}
		\item[($ii$)] If $X$ has independent and stationary increments,
		\begin{equation}\label{incrementoIntegraleProcessoStocasticoIncrementiIndipendentiStazionario}
			\mathbb{E}e^{i\gamma \big(I(t_2)-I(t_1)\big)} = \mathbb{E}e^{i\gamma (t_2-t_1)X(t_1)}\mathbb{E}e^{i\gamma I(t_2-t_1)}.
		\end{equation}
	\end{itemize}
\end{lemma}
See Apeendix \ref{appendiceDimostrazioneIncrementiIntegrale} for the proof of Lemma \ref{lemmaIncrementiIntegraleProcesso}.





\subsection{Poisson case}

For the sake of completeness we recall the following result which follows by classical arguments.

\begin{lemma}\label{lemmaLimiteSerieRappresentanteIntegrale}
	Let $g:\mathbb{N}^2\longrightarrow \mathbb{R}$ and $k:\mathbb{R}\longrightarrow\mathbb{R}$ such that for $\varepsilon >0\ \exists\ N_\varepsilon$ s.t. $|g(h,n) - k(h/n)|<\varepsilon,\ n\ge N_\varepsilon,\ h=1,\dots, n$. Let $f:\mathbb{R}\longrightarrow\mathbb{R}$ s.t. $\int_0^1 |f(x)|\dif x<\infty$ and $\exists\ \int_0^1 f(x)k(x)\dif x$. Then, 
	\begin{equation*}
		\lim_{n\rightarrow\infty} \sum_{h=1}^n \frac{1}{n}f\Big(\frac{h}{n}\Big) g(h,n) = \int_0^1f(x)k(x)\dif x.
	\end{equation*}
\end{lemma}

Now we prove that the fractional integral of a compound Poisson process (up to a constant) is a compound Poisson random variable.

\begin{theorem}\label{teoremaIntegraleProcessoCompostoPoisson}
Under Condition II, let $M\sim CP(\lambda, X)$ and $\alpha\ge0$. For $t>0$, set $\Lambda_t$ a r.v. such that
\begin{equation}\label{distribuzioneVariabileAleatoriaLambda}
	P\{\Lambda_t\in \dif s\} = \frac{\lambda(s)\dif s}{\int_0^t \lambda(u)\dif u}, \ \ \ s\in[0,t].
\end{equation}
Then,
\begin{equation}\label{integraleProcessoCompostoPoisson}
	\int_0^t (t-s)^\alpha M(s)\dif s \stackrel{d}{=} \sum_{k=1}^{N(t)} X_k\frac{(t-\Lambda_{t,k})^{\alpha+1}}{\alpha+1},\ \ \ t\ge0,
\end{equation}	
where $\Lambda_{t,1}, \Lambda_{t,2},\dots$ are independent copies of $\Lambda_t$, independent of $N$ and of $X_1,X_2,\dots$.
\end{theorem}

\begin{proof}
	First, we observe that the characteristic function of the right hand-side of \eqref{integraleProcessoCompostoPoisson} is, for $\gamma \in\mathbb{R}$,
	\begin{equation}
		\mathbb{E}\exp\Bigg( i\gamma \sum_{k=1}^{N(t)} X_k\frac{(t-\Lambda_{t,k})^{\alpha+1}}{\alpha+1} \Bigg)= e^{-\int_0^t \lambda(s)\dif s} \exp\Bigg( \int_0^t \lambda(s) \mathbb{E} e^{i\gamma X \frac{(t-s)^{\alpha+1}}{ \alpha+1} }\dif s \Bigg).
	\end{equation}
	Now, we study the characteristic function of the left hand-side of \eqref{integraleProcessoCompostoPoisson}.
\begin{align}
\int_0^t (t-s)^\alpha M(s)\dif s & = 	\lim_{n\rightarrow \infty} \sum_{k=1}^n \bigg(t-\frac{kt}{n}\bigg)^\alpha M\bigg(\frac{kt}{n}\bigg)\frac{t}{n}\nonumber \\
& =  \lim_{n\rightarrow \infty} \sum_{h=1}^n \Bigg(M\bigg(\frac{ht}{n}\bigg) - M\bigg(\frac{(h-1)t}{n}\bigg) \Bigg) \sum_{k = h}^n \bigg(\frac{t}{n}\bigg)^{\alpha+1}(n-k)^\alpha \nonumber \\
&= \lim_{n\rightarrow \infty} \sum_{h=1}^n \Bigg(M\bigg(\frac{ht}{n}\bigg) - M\bigg(\frac{(h-1)t}{n}\bigg) \Bigg) \sum_{j = 0}^{n-h} \bigg(\frac{t}{n}\bigg)^{\alpha+1} j^\alpha .\label{primaRappresentazioneIntegraleFrazinarioCompound}
\end{align}

Note that for $n$ great enough, then, for $ h=1,\dots,n$ and $\gamma \in\mathbb{R}$,
\begin{align}\label{funzioneCaratteristicaIncrementoCompoundPoisson}
	\mathbb{E} \exp\Bigg(i\gamma \bigg[ M\bigg(\frac{ht}{n}\bigg) - M\bigg(\frac{(h-1)t}{n}\bigg)\bigg] \Bigg) =1-\lambda\bigg(\frac{(h-1)t}{n}\bigg)\frac{t}{n}\Bigg(1-\mathbb{E}e^{i\gamma X}\Bigg)+o\bigg(\frac{t}{n}\bigg).
\end{align}
Then, the characteristic function of \eqref{primaRappresentazioneIntegraleFrazinarioCompound} reads
\begin{align}
	\mathbb{E}\exp\bigg(&i\gamma \int_0^t (t-s)^\alpha M(s)\dif s\bigg)\nonumber \\
	&= \lim_{n\rightarrow\infty}\prod_{h=1}^n\mathbb{E} \exp\Bigg(i\gamma \bigg[M\bigg(\frac{ht}{n}\bigg) - M\bigg(\frac{(h-1)t}{n}\bigg) \bigg] \sum_{j = 0}^{n-h} \bigg(\frac{t}{n}\bigg)^{\alpha+1} j^\alpha \Bigg)\nonumber\\
	& = \lim_{n\rightarrow\infty}\prod_{h=1}^n \Bigg( 1-\lambda\bigg(\frac{(h-1)t}{n}\bigg)\frac{t}{n}\bigg[1-\mathbb{E}\exp\bigg( i\gamma X \sum_{j = 0}^{n-h} \bigg(\frac{t}{n}\bigg)^{\alpha+1} j^\alpha \bigg)\bigg] \Bigg)\nonumber\\
	& = \exp\Bigg( -\lim_{n\rightarrow\infty} \sum_{h=1}^n\lambda\bigg(\frac{(h-1)t}{n}\bigg)\frac{t}{n}\bigg[1-\mathbb{E}\exp\bigg( i\gamma X \sum_{j = 0}^{n-h} \bigg(\frac{t}{n}\bigg)^{\alpha+1} j^\alpha \bigg)\bigg] \Bigg) \nonumber\\ % COMMENTA
	& = e^{-\int_0^t \lambda(s)\dif s} \exp\Bigg(\lim_{n\rightarrow\infty} \sum_{h=1}^n \lambda\bigg(\frac{(h-1)t}{n}\bigg)\frac{t}{n} \mathbb{E}\exp\bigg( i\gamma X \sum_{j = 0}^{n-h} \bigg(\frac{t}{n}\bigg)^{\alpha+1} j^\alpha \bigg) \Bigg) ,\nonumber
\end{align}
where in the first step we used \eqref{primaRappresentazioneIntegraleFrazinarioCompound}, the dominated convergence theorem to exchange limit and expectation and the independence of the increments of $M$.  In the second step we used \eqref{funzioneCaratteristicaIncrementoCompoundPoisson}. In the third step we used that if $\lim_{n\rightarrow\infty} \sum_{k=1}^n \frac{a_{k,n}}{n}<\infty$ and $\sum_{k=1}^n \frac{a^2_{k,n}}{n^2}=0$ (which can be easily checked), then
$$\lim_{n\rightarrow\infty} \prod_{k=1}^n\Big(1+\frac{a_{k,n}}{n}\Big) = \exp\Bigg( \lim_{n\rightarrow\infty}\sum_{k=1}^n \frac{a_{k,n}}{n}\Bigg).$$

Now, it remains to prove that
\begin{equation}
	\lim_{n\rightarrow\infty} \sum_{h=1}^n \lambda\bigg(\frac{(h-1)t}{n}\bigg)\frac{t}{n} \mathbb{E}\exp\bigg( i\gamma X \sum_{j = 0}^{n-h} \bigg(\frac{t}{n}\bigg)^{\alpha+1} j^\alpha \bigg) = \int_0^t \lambda(s) \mathbb{E} e^{i\gamma X \frac{(t-s)^\alpha}{ \alpha+1}  } \dif s
	\end{equation}

In light of Lemma \ref{lemmaLimiteSerieRappresentanteIntegrale}, it sufficies to show that $ \sum_{j = 0}^{n-h} \big(\frac{t}{n}\big)^{\alpha+1} j^\alpha \sim \frac{(t-ht/n)^{\alpha+1}}{\alpha+1}$ as $n$ increases. Indeed,

\begin{align}
	\sum_{j=0}^{n-h} \frac{j^\alpha}{n^{\alpha+1}}  & = \frac{1}{n}\sum_{j=0}^{n-h}\sum_{k\ge0} \frac{\Gamma(\alpha+1)}{k!\Gamma(\alpha+1-k)}  \Big(\frac{j}{n}-1\Big)^k \nonumber \\
	&= \frac{1}{n}\sum_{j=0}^{n-h}\sum_{k\ge0} \frac{\Gamma(\alpha+1)}{k!\Gamma(\alpha+1-k)}  \sum_{i=0}^k \binom{k}{i} (-1)^{k-i} \Big(\frac{j}{n}\Big)^i \nonumber\\ % COMMENTA
	& = \frac{1}{n}\sum_{k\ge0} \frac{\Gamma(\alpha+1)}{k!\Gamma(\alpha+1-k)}  \sum_{i=0}^k \binom{k}{i} (-1)^{k-i} \sum_{j=0}^{n-h}\Big(\frac{j}{n}\Big)^i \nonumber\\
	&= \sum_{k\ge0} \frac{\Gamma(\alpha+1)}{k!\Gamma(\alpha+1-k)}  \sum_{i=0}^k \binom{k}{i} (-1)^{k-i} \frac{1}{i+1}\sum_{j=0}^i \binom{i+1}{j} (-1)^j B_j \frac{(n-h)^{i+1-j}}{n^{i+1}}, \label{espansioneSommatoriaPotenzaReale} %\label{passaggioFormulaFaulhaber}
\end{align}
where in the last step we used the Faulhaber's formula $\sum_{j=1}^m j^i = \sum_{j=0}^i \binom{i+1}{j} (-1)^j B_j m^{i+1-j} / (i+1)$, $i\ge1$, where $B_j$ are the Bernoulli numbers with $B_0 = 1$. Then, for $n$ big enough the last sum in \eqref{espansioneSommatoriaPotenzaReale} reduces to the term with $j=0$, i.e.
\begin{align}
	\sum_{j=0}^{n-h} \frac{j^\alpha}{n^{\alpha+1}}  & \sim \sum_{k\ge0} \frac{\Gamma(\alpha+1)}{k!\Gamma(\alpha+1-k)}  \sum_{i=0}^k \binom{k}{i} (-1)^{k-i} \frac{1}{i+1} \bigg(1-\frac{h}{n}\bigg)^{i+1}  \nonumber\\
	& = \frac{1}{\alpha+1}\sum_{k\ge0} \frac{\Gamma(\alpha+2)}{(k+1)!\Gamma(\alpha+1-k)} (-1)^{k+1} \sum_{i=0}^k \binom{k+1}{i+1} (-1)^{i+1} \bigg(1-\frac{h}{n}\bigg)^{i+1}  \nonumber\\
	& =  \frac{1}{\alpha+1}\sum_{k\ge0} \frac{\Gamma(\alpha+2)}{(k+1)!\Gamma(\alpha+1-k)} (-1)^{k+1} \Bigg( \bigg(\frac{h}{n}\bigg)^{k+1}- 1\Bigg) \nonumber\\
	& = \frac{1}{\alpha+1} \bigg(1-\frac{h}{n}\bigg)^{\alpha +1},\nonumber
\end{align}
which completes the proof.
\end{proof}


\begin{remark}[Iterated integrals]
	In light of the Cauchy integral formula, the iterated integral of $M\sim CP(\lambda, X)$ is, with $t\ge0$ and $n\in\mathbb{N}$,
	\begin{align*}
		\int_0^t\dif s_1\int_0^{s_1} \dif s_2\cdots \int_0^{s_{n-1}}M(s_n)\dif s_n &= \frac{1}{(n-1)!}\int_0^t M(s)(t-s)^{n-1} \dif s \nonumber\\
		& \stackrel{d}{=} \frac{t^n}{n!} \sum_{k=1}^{N(t)} X_k \Big(1-\frac{\Lambda_{t,k}}{t}\Big)^{n} \stackrel{p}{\longrightarrow}0, \ \ \ n\rightarrow\infty,
	\end{align*}
	where the limit follows by assuming that $N$ does not explode (see Condition II) since $0<1-\Lambda_{t,k}/t<1$ a.s. \hfill $\diamond$
\end{remark}

\begin{remark}[Homogeneous case]\label{remarkIntegraleCompoundOmogeneo}
	If $\lambda(t) = \lambda ,\ t\ge0$, the formula \eqref{integraleProcessoCompostoPoisson} simplifies. Indeed, $\Lambda_t \sim Unif(0,t)$ and $t-\Lambda_t \sim Unif(0,t)$. Thus, denoting by $U_k\sim Unif(0,1)$ i.i.d. r.v.s independent of $N$ and $X_1,\dots$,
	\begin{equation}\label{integraleFrazionarioCompoundOmogeneo}
		\int_0^t M(s)(t-s)^\alpha \dif s \stackrel{d}{=} \frac{t^{\alpha+1}}{\alpha+1} \sum_{k=1}^{N(t)} X_k U_k^{\alpha+1},\ \ \ t\ge0,\ \alpha\ge0.
	\end{equation}
	Note that the case of $\alpha = 0$ already appeared in \cite{CO202}.
	
	We readily derive the following limit result. Without any loss of generality, let $\mathbb{E}X =0$ and $\mathbb{E}X^2 <\infty$, then as $t\longrightarrow \infty$,
	\begin{align*}
		\lim_{t\rightarrow\infty} \frac{1}{t^{\beta}} \int_0^t M(s) \dif s \stackrel{d}{=} \frac{1}{t^{\beta-1}} \sum_{k=1}^{N(t)} X_k U_k 	
		 \longrightarrow \begin{cases}
			\begin{array}{l l}
				\displaystyle \stackrel{p}{\longrightarrow} 0, & \beta = 2,\\
				\displaystyle \stackrel{d}{\longrightarrow} Z\sim Normal(0, \lambda \mathbb{E}X^2 / 3\big), & \beta = 3/2,
			\end{array}
		\end{cases}
	\end{align*}
	where we used \eqref{integraleFrazionarioCompoundOmogeneo} and the limit theorems for the sum of a random number of i.i.d. terms. \hfill$\diamond$
\end{remark}

Now, we point out an interesting connection between Theorem \ref{teoremaIntegraleProcessoCompostoPoisson} and formula \eqref{formulazioneGeneraleIntegraleProcessoCompound} with $f(s) = (t-s)^\alpha$. Indeed, the latter turns into
\begin{equation}\label{rappresentazioneQuasiCertantegraleCompoundPoisson}
	\int_0^t M(s)(t-s)^{\alpha}\dif s = \sum_{k=1}^{N(t)} X_k \frac{(t-S_k)^\alpha}{\alpha+1},\ \ \ t\ge0,
\end{equation} 
which resembles \eqref{integraleProcessoCompostoPoisson}, where the interarrival times $S_k$ are replaced by the r.v.s $\Lambda_{t,k}$. The following results explore this idea and permit us to generalize Theorem \ref{teoremaIntegraleProcessoCompostoPoisson}. We also point out that formula \eqref{rappresentazioneQuasiCertantegraleCompoundPoisson} is a particular shot noise process, whose basic formulation is $S(t) = \sum_{k=1}^{N(t)} X_k f(t-S_k)$ with $f$ usually decreasing (see for instance \cite{BCS2026} and references therein).

\begin{proposition}\label{proposizioneDistribuzioneCongiuntaIntertempiPoisson}
	Let $N\sim PP(\lambda)$. Let $t>0$, $\Lambda_{t,1},\dots$ i.i.d r.v.s distributed as in \eqref{distribuzioneVariabileAleatoriaLambda}. Then, for $n\in\mathbb{N}$ and $0\le s_1<s_2<\dots<s_n\le t$,
	\begin{equation}
		P\{S_1\in\dif s_1,\dots,S_n\in \dif s_n\,|\,N(t) = n\} = P\{\Lambda_{t,(1)}\in\dif s_1, \dots, \Lambda_{t, (n)}\in\dif s_n\},
	\end{equation}
	with $\Lambda_{t,(1)},\dots,\Lambda_{t,(n)}$ the (increasingly) ordered $\Lambda_{t,1},\dots,\Lambda_{t,n}$.
\end{proposition}

\begin{proof}
	\begin{align*}
		P\{S_1\in \dif s_1,&\dots,S_n\in \dif s_n\,|\,N(t) = n\} \\
		&= \frac{P\{N[o,s_1) = 0,N[s_1,s_1+\dif s_1) = 1,\dots, N[s_n,s_n+\dif s_n) = 1,N[s_n+\dif s_n,t) = 0\}}{P\{N(t) = n\}} \\
		&= \frac{n!}{\big(\int_0^t\lambda(s)\dif s\big)^n} \prod_{k=1}^n\lambda(s_k)\dif s_k \\
		&= 	P\{\Lambda_{t,(1)}\in\dif s_1, \dots, \Lambda_{t, (n)}\in\dif s_n\}.
	\end{align*}
\end{proof}

\begin{proposition}\label{proposizioneVariabiliOrdinateFunzioniCommutative}
	Let $n\in\mathbb{N}$ and $t>0$. Let $\Lambda_1,\dots,\Lambda_n$ i.i.d. r.v.s taking values in $[0,t]$ a.s., $X_1,\dots,X_n$ independent copies of a real r.v. $X$ and $f:\mathcal{R}\times [0,t] \longrightarrow \mathbb{R}$ measurable. Then,
	\begin{equation}
		\sum_{k=1}^n f\Big(X_k, \Lambda_{(k)}\Big) \stackrel{d}{=} \sum_{k=1}^n f\big( X_k,\Lambda_k \big).
	\end{equation}
\end{proposition}

Note that the underlying reason is that the sum is commutative and, thanks to the independence and identical distribution of the $X_k$, the order of the terms does not matter. Indeed, we could replace the summation with any commutative (with respect to the inputs) function $g$, obtaining 
$$ g\Bigg(f\Big(X_1, \Lambda_{(1)}\Big),\dots, f\Big(X_n, \Lambda_{(n)}\Big) \Bigg) \stackrel{d}{=} g\Bigg(f\Big(X_1, \Lambda_{1}\Big),\dots, f\Big(X_n, \Lambda_{n}\Big) \Bigg). $$

\begin{proof}
	
	Let $\gamma\in\mathbb{R}$, and $\Lambda$ distributed as $\Lambda_k$,
	\begin{align}
		\mathbb{E} \exp\Bigg(i\gamma\sum_{k=1}^n f(X_k,\Lambda_k) \Bigg) &= \left( \mathbb{E} e^{i\gamma f(X,\Lambda) }\right)^n \nonumber\\
		& =  \Bigg( \int_0^t \mathbb{E} e^{i\gamma f(X,\lambda)}P\{\Lambda \in\dif \lambda\} \Bigg)^n \nonumber\\
		& = n! \int_0^t \int_{\lambda_1}^t \cdots \int_{\lambda_{n-1}}^t \prod_{k=1}^n \Bigg(\mathbb{E} e^{i\gamma f(X,\lambda_k)}P\{\Lambda \in\dif \lambda_k\} \Bigg), \label{passaggioUsoFormulaPotenzaIntegrale} \\
		&=  n! \int_0^t \int_{\lambda_1}^t \cdots \int_{\lambda_{n-1}}^t \mathbb{E} \exp\Big(i\gamma \sum_{k=1}^n f(X_k,\lambda_k)  \Big) \prod_{k=1}^n  P\{\Lambda \in\dif \lambda_k\} \nonumber \\
		& = \mathbb{E} \exp\Bigg(i\gamma\sum_{k=1}^n f(X_k,\Lambda_{(k)}) \Bigg),\nonumber
	\end{align}
	%\begin{lemma}	
	where in the second last equality we used that $X_1,\dots,X_n$ are independent copies of $X$ and in \eqref{passaggioUsoFormulaPotenzaIntegrale} we used that, if $f:[0,\infty)\longrightarrow \mathbb{C}$, then, for $t>0$ and $n\in\mathbb{N}$,
	\begin{equation}
		\Biggl(\int_0^t f(x)\dif x \Bigg)^n = n! \int_0^t f(x_1) \dif x_1\int_{x_1}^t f(x_2)\dif x_2 \cdots\int_{x_{n-1}}^t f(x_n)\dif x_n,
	\end{equation}
	which can be proved by induction.
	%\end{lemma}
	\end{proof}


Now, equation \eqref{formulazioneGeneraleIntegraleProcessoCompound}, Proposition \ref{proposizioneDistribuzioneCongiuntaIntertempiPoisson} and Proposition \ref{proposizioneVariabiliOrdinateFunzioniCommutative} yield the following theorem, meaning that for fixed $t\in[0, T]$ the integral \eqref{formulazioneGeneraleIntegraleProcessoCompound} of a compound Poisson process is a compoud Poisson random variable.


\begin{theorem}\label{teoremaIntegraleProcessoCompostoPoissonGenerale}
	Let $T>0$ and measurable $f\in \mathcal{L}^1[0,T]$. Under the hypothesis of Theorem \ref{teoremaIntegraleProcessoCompostoPoisson}, 
	\begin{equation}\label{integraleGeneraleProcessoCompostoPoisson}
		\int_0^t  M(s)f(s)\dif s \stackrel{d}{=} \sum_{k=1}^{N(t)} X_k\int_{\Lambda_{t,k}}^t f(s)\dif s,\ \ \ 0\le t \le T.
	\end{equation}	
\end{theorem}

We point out that the previous result extends to the shot noise processes as well, meaning that $S(t) = \sum_{k=1}^{N(t)} X_k f(t-S_k) \stackrel{d}{=} \sum_{k=1}^{N(t)} X_k f(t-\Lambda_{t,k})$. Furthermore, in view of point ($i$) of Remark \ref{remarkUguaglianzaSuSpazioFunzioniContinueCompoundDiCompound} the compound of a compound Poisson process still is a compound Poisson process (since the equality in distribution over $\mathcal{D}[0,T]$) and therefore Theorem \ref{teoremaIntegraleProcessoCompostoPoissonGenerale} suitably holds for the integral of the "iterated" compound Poisson processes.

\begin{remark}
	We show that, in general, both relationship \eqref{integraleProcessoCompostoPoisson} and \eqref{integraleGeneraleProcessoCompostoPoisson} do not hold on $\mathcal{D}[0,T]$. It sufficies to consider the case of homogeneous $M\sim CP(\lambda, X)$ and $f\equiv 1$. In light of Remark \ref{remarkIntegraleCompoundOmogeneo} and formula  \eqref{incrementoIntegraleProcessoStocasticoIncrementiIndipendentiStazionario}, with $0\le t_1<t_2\le T$ and $\gamma \in\mathbb{R}$, the increment of the integral is such that
	\begin{equation}\label{incrementoVeroIntegralePoissonCompostoOmogeneo}
		\mathbb{E}e^{i\gamma \int_{t_1}^{t_2} M(s)\dif s} = \exp\Big(-\lambda t_2 + \lambda t_1 \mathbb{E} e^{i\gamma (t_2-t_1)X} +\lambda(t_2 -t_1)\mathbb{E}e^{i\gamma (t_2-t_1) XU}  \Big),
	\end{equation}
	with $U\sim Unif(0,1)$ being independent of $X$. On the other hand, denoting by $Z(t) = t\sum_{k=1}^{N(t)} X_k U_k, \ 0\le t\le T,$ as in \eqref{integraleFrazionarioCompoundOmogeneo}, we have 
	\begin{align}
		\mathbb{E}e^{i\gamma \big(Z(t_2)-Z(t_1)\big) } & = \mathbb{E} \exp\Bigg(i\gamma (t_2-t_1)\sum_{k=1}^{N(t_1)}X_kU_k + i\gamma t_2 \sum_{k=1}^{N(t_2)-N(t_1)} X_kU_k  \Bigg) \nonumber\\
		& = \exp\Bigg( -\lambda t_2 +\lambda t_1\Big( \mathbb{E}e^{i\gamma (t_2-t_1)XU}\big)+\lambda(t_2 -t_1)\mathbb{E}e^{i\gamma t_2 XU} \Bigg), \nonumber %\label{incrementoRappresentazioneIntegralePoissonCompostoOmogeneo}
	\end{align}
	which in general differs from \eqref{incrementoVeroIntegralePoissonCompostoOmogeneo}. \hfill$\diamond$
\end{remark}


\begin{remark}[Limit behaviour]
	Inspired by the Kac's hydrodynamic limit conditions, Theorem \ref{teoremaIntegraleProcessoCompostoPoissonGenerale} permits us to derive the limit result when the rate function of the underlying Poisson process goes to infinity. Let $\beta >0$ and $M\sim CP(\lambda, X)$ with integrable $\lambda$ such that $\lambda(t) = \beta \xi(t),\ t\ge0$. Now, the r.v. $\Lambda_{t}$, having density \eqref{distribuzioneVariabileAleatoriaLambda}, does not depend on $\beta$ (since its distribution does not). Furthermore, $N(t)/\beta \stackrel{p,L^1}{\longrightarrow} \xi(t)\ \forall\ t$ (see for instance Corollary 2.1 of \cite{CO202}). Hence, we have the following limit result, for $\mathbb{E}X = 0$ and $\mathbb{E}X^2<\infty$,
	\begin{equation}\label{comportamentoLimiteTipoKAc}
		\frac{1}{\beta^\alpha} \int_0^t M(s)f(s)\dif s 
	 \longrightarrow \begin{cases}
			\begin{array}{l l}
				\displaystyle \stackrel{p}{\longrightarrow} 0, & \alpha = 1,\\
				\displaystyle \stackrel{d}{\longrightarrow} Z\sim Normal\Bigg(0, \xi(t) \mathbb{E}X^2 \mathbb{E}\bigg(\int_{\Lambda_t}^t f(s)\dif s\bigg)^2\Bigg), & \alpha = 1/2,
			\end{array}
		\end{cases}
	\end{equation}
	as $\beta\longrightarrow\infty$.
	\\
	
	As an example, inspired by the paper \cite{MS2025} (also see references therein), we consider $\lambda(t) = \beta /t^{\alpha}$ with $\alpha \in(0,1)$ and $f(s) = s^r$ with $r>0$. Thus, $\mathbb{E}\Lambda_t^u = t^u \frac{1-\alpha}{1-\alpha+u}, \ u\ge 1$. In this case, 
	\begin{align*}
		\mathbb{E} \bigg(\int_{\Lambda_t}^t s^r\dif s\bigg)^2 & = \frac{t^{2(r+1)}-2t^{r+1}\mathbb{E}\Lambda_t^{r+1}+\mathbb{E}\Lambda_t^{2(r+1)}}{(r+1)^2} \\
		&=\frac{t^{2(r+1)}}{(r+1)^2} \bigg( 1-\frac{2(1-\alpha)}{1-\alpha+r+1}+\frac{1-\alpha}{1-\alpha+2(r+1)}\bigg).
		\end{align*}
		Finally, by means of \eqref{comportamentoLimiteTipoKAc},
$$\frac{1}{\sqrt{\beta}} \int_0^t M(s)s^r\dif s \stackrel{d}{\longrightarrow} Z\sim Normal\Bigg(0, \mathbb{E}X^2 \frac{t^{2(r+1)-\alpha}}{(r+1)^2} \bigg( 1-\frac{2(1-\alpha)}{2-\alpha+r}+\frac{1-\alpha}{3-\alpha+2r)}\bigg)\Bigg),  $$
	as $\beta \longrightarrow\infty$.
	\hfill$\diamond$
\end{remark}

\subsubsection{Governing partial difference-differential equation}

Under Condition II, let $T>0$ and $Y(t) = \int_0^t g\big(M(s)\big) f(s)\dif s,\ 0\le t\le T,$ where $M\sim CP(\lambda, X)$, $f\in \mathcal{L}^1\big([0,T]\big)$ and measurable $g: \mathbb{R}\longrightarrow \mathbb{R}$, $g\in\mathcal{L}^1(\mathbb{R})$. Following the line of \cite{GMn1971, Mn1970, P1966, P1968}, we study the equation governing the probability law of the vector process $(M,Y)$, $p(t,\dif x, y) \dif y = P\{M(t)\in\dif x, Y(t)\in\dif y\}$ (where the $\dif x$ term must be interpreted as in \eqref{compoundPoissonInfinitesimale}). For $t\ge0,\ x,y \in\mathbb{R},$

\begin{align}
	p(t+\dif t,\dif x, y) &= P\{M(t)+M[t,t+\dif t) \in \dif x, Y(t)\in\dif y - \dif t g\big(M(t)\big)f(t) \} \nonumber\\
	& = \int_\mathbb{R} P\{M(t)\in \dif z, Y(t)\in \dif y-g(z)f(t)\}P\{M[t,t+\dif t)\in\dif x-z\} \nonumber\\
	& = \int_\mathbb{R} p\big(t, \dif z, y-\dif t g(z)f(t)\big)\lambda(t)\dif t F(\dif x -z)\nonumber \\
	&\ \ \ +p\big(t, \dif x, y-\dif t g(x) f(t)\big) \big(1-\lambda(t)\dif t\big)+ o(\dif t),\nonumber
\end{align}
and by expanding and keeping the terms of order $\dif t$ or lower,
\begin{align}
	&p(t,\dif x, y) + \frac{\partial}{\partial t}p(t,\dif x,y)\dif t\label{formulaEquazioneDifferenzialePrimoOrdine}\\
	&=  \lambda(t)\dif t\int_\mathbb{R} p(t,\dif x-z, y)F(\dif z) - \big(1-\lambda(t)\dif t)\big)p(t,\dif x,y) -  g(x)f(t) \frac{\partial}{\partial y}p(t,\dif x, y) +o(\dif t).\nonumber
\end{align}
By dividing by $\dif t$ and taking the limit $\dif t\longrightarrow 0$, equation \eqref{formulaEquazioneDifferenzialePrimoOrdine} turns into
	\begin{align}
	\frac{\partial}{\partial t}p(t, \dif x, y)+ g(x)f(t)& \frac{\partial}{\partial y}p(t, \dif x, y) = \lambda(t) \Big(I-C_F^{(x)}\Big) p(t, \dif x, y) , \ \ \ t\ge0,\ x,y \in\mathbb{R},
\end{align}
where $C_{F}^{(x)} p(t,\dif x, y) =\int_\mathbb{R} p(t,\dif x - z, y) F(\dif z)$.
\\

By suitably adapting the above steps, we derive the following theorem concerning the differential equation governing the $(n+1)$-th dimensional vector describing the compound process and its $n$ iterated integrals.

\begin{theorem}
	Let $T>0$ and $n\in\mathbb{N}$. Let $M\sim CP(\lambda, F)$, measurable real functions $g,f_1,\dots,f_n\in\mathcal{L}^1[0,T]$. For $0\le t\le T$, define
	$$Y_1(t) = \int_0^t g\big(M(s)\big) f_1(s) \dif s,\ \ \ Y_k(t) = \int_0^t Y_{k-1}(s)f_k(s)\dif s,\ k=2,\dots, n. $$
	Then, the probability law $p(t,\dif x, y)\dif y = P\{M(t)\in\dif x, Y_1(t)\in \dif y_1,\dots, Y_n(t)\in \dif y_n\}$ satisfies
	\begin{align}
		\frac{\partial}{\partial t}p(t, \dif x, y)+ g(x)f_1(t)& \frac{\partial}{\partial y_1}p(t, \dif x, y) + \sum_{k=2}^n y_{k-1} f_k(t)\frac{\partial}{\partial y_{k}}p(t, \dif x, y) \nonumber\\
		& = \lambda(t) \Big(I-C_F^{(x)}\Big) p(t, \dif x, y) , \ \ \ t\in [0,T],\ x,y_1,\dots,y_n \in\mathbb{R}, \label{equazioneCompoundPoissonIntegraliIterati}
	\end{align}	
	subject to $p(0,\dif x, y) = \delta_0(\dif x, y)$.
\end{theorem}

Note that the operator on the right-hand side depends on $x$ and $t$ only and it has the same structure of the one in \eqref{equazioneOperatorialeCompoundPoisson}. 

\begin{remark}[Fractional equations and time-changing]
	
We point out that in the homogeneous case, by assuming $g(x)=x,x\in\mathbb{R}, f_1\equiv\dots\equiv f_n\equiv1$ and $T=\infty$, it is easy to derive some results concerning the composition with a subordinator or its inverse. In particular, by applying a Bernstein function $f$ to the right-hand side of \eqref{equazioneCompoundPoissonIntegraliIterati}, using the arguments of point ($ii$) of Theorem \ref{teoremaComposizioneCompoundPoissonConSubordinatore} we readily obtain that the solution is still in the class of the compound Poisson processes and it concerns the process $M\circ \mathcal{H}_f$ and its iterated integrals, where $\mathcal{H}_f$ is a Bernstein subordinator.
\\

On the other hand, following the arguments of Proposition \ref{proposizioneComposizioneInversoSubordinatore}, the probability law $p_f$ of the time-changed vector $(M\circ L_f, Y_1\circ L_f,\dots, Y_n\circ L_f)$ (i.e. it is not the vector concerning $M\circ L_f$ and its iterated integrals), where $L_f$ is an inverse Bernstein subordinator, satisfies the fractional problem, with $t\ge0,\ x,y_1,\dots,y_n\in \mathbb{R}$,
\begin{equation}\label{equazioneOperatorialeCompoundIntegraliIterati}
	\mathcal{D}_t^f p_f(t,\dif x,y) = -\lambda\Big(I-C_F^{(x)}\Big) p_f(t,\dif x,y) - x \frac{\partial}{\partial y_1}p_f(t, \dif x, y) - \sum_{k=2}^n y_{k-1} \frac{\partial}{\partial y_{k}}p_f(t, \dif x, y),
\end{equation}
subject to $p_f(0,\dif x,y) = \delta_0(\dif x,y)$and under the hypotheses $\mathcal{D}_t^f p_f(t,\dif x,y) = \int_0^\infty p(s,\dif x,y) \mathcal{D}_t^f l_f(t,s)\dif s$
and $\int_0^\infty O_{x,y}p(s, \dif x, y)l_f(t,s)\dif s = O_{x,y} p_f(t,\dif x, y)$, with $O_{x,y}$ being the space operator on the right-hand side of \eqref{equazioneOperatorialeCompoundIntegraliIterati}.
\hfill$\diamond$
\end{remark}

\subsubsection{Homogeneous Skellam-type version}

We consider the homogeneous Skellam-type process described in Example \ref{esempioSkellam}, i.e. $M = \sum_{i \in\mathcal{I} } i N_i$, with finite $\mathcal{I}$ and $N_i\sim PP(\lambda_i), \ i\in\mathcal{I}\subset \mathbb{R}\setminus \{0\}$, being independent homogeneous Poisson processes. 

\begin{proposition}
	Let $M$ be a homogeneous Skellam-type process defined above. Then, for $x>0$,
\begin{align}
	P\bigg\{&\int_0^t M(s)\dif s\in \dif x\bigg\} \nonumber\\
	&=  e^{-\lambda t} \sum_{n\ge0} \frac{1}{n!} \sum_{\substack{k_1,\dots,k_I=0\\k_1+\dots+k_I=n}}^n \binom{n}{k_1,\dots, k_I} \prod_{i\in\mathcal{I}} \Bigg( \bigg(\frac{\lambda_i}{i}\bigg)^{k_{(i)}}  \sum_{h_{(i)} = 0}^{k_{(i)}} \binom{k_{(i)}}{h_{(i)}} (-1)^{h_{(i)}} \Bigg)\nonumber \\
	&\ \ \ \times\mathds{1}_{(t\sum_{j\in\mathcal{I}} j h_{(j)} ,\infty)}(x) \frac{\big(x-t\sum_{j\in\mathcal{I}} j h_{(j)}\big)^{n-1}}{(n-1)!}\dif x, \label{probabilitaEsplicitaIntegraleSkemllamGeneralizzato}
\end{align}
where $I = |\mathcal{I}|<\infty$ and, for $i\in\mathcal{I}$, $(i)$ denotes the index of $i$ in the increasingly ordered version of $\mathcal{I}$ (thus, for instance $(\min\{\mathcal{I}\}) = 1$ and $(\max\{\mathcal{I}\}) = I$).
\end{proposition}

\begin{proof}
Write $M$ in the following way as a compound Poisson process, $M\sim CP( \lambda = \sum_{i\in\mathcal{I}} \lambda_i, Y)$, with $P\{Y = i\} = \lambda_i/\lambda$. Now, Remark \ref{remarkIntegraleCompoundOmogeneo} permits us to easily write the moment generating function of the integral of $M$. Let us also assume  $U\sim Unif(0,1)$ independent of $Y$ and $t>0,\ \mu\in\mathbb{R}$ ($M$ can take real values with positive probability, so we consider the bilateral Laplace transform), then,
\begin{align}
	\mathbb{E}&\exp\bigg( -\mu \int_0^t M(s)\dif s\bigg) \nonumber\\
	& = \exp\bigg(-\lambda t\big(1-\mathbb{E}e^{-\mu t YU}\big)\bigg)\nonumber \\
	& = \exp\Bigg(-\lambda t \bigg(1- \sum_{i\in\mathcal{I}}\frac{\lambda_i}{\lambda}\, \frac{1-e^{-\mu t i}}{\mu t i}\bigg) \Bigg)\nonumber\\
	%& = e^{-\lambda t}\sum_{n\ge0} \frac{1}{n!} \bigg( \sum_{i\in\mathcal{I}} \lambda_i t \text{sgn}(i) \frac{1-e^{-\mu t i}}{\mu}\bigg)^n\nonumber\\
	& = e^{-\lambda t} \sum_{n\ge0} \frac{\mu^{-n}}{n!}  \Bigg( \sum_{i\in\mathcal{I}} \frac{\lambda_i}{i} \Big(1-e^{-\mu t i}\Big)\Bigg)^n\nonumber\\
	& = e^{-\lambda t} \sum_{n\ge0}\frac{\mu^{-n}}{n!}  \sum_{\substack{k_1,\dots,k_I=0\\k_1+\dots+k_I=n}}^n \binom{n}{k_1,\dots, k_I} \prod_{i\in\mathcal{I}} \bigg(\frac{\lambda_i}{i}\bigg)^{k_{(i)}}  \Big( 1-e^{-\mu t i}\Big)^{k_{(i)}}\nonumber\\
	& = e^{-\lambda t} \sum_{n\ge0} \frac{\mu^{-n}}{n!}  \sum_{\substack{k_1,\dots,k_I=0\\k_1+\dots+k_I=n}}^n \binom{n}{k_1,\dots, k_I} \prod_{i\in\mathcal{I}} \bigg(\frac{\lambda_i}{i}\bigg)^{k_{(i)}}  \sum_{h_{(i)} = 0}^{k_{(i)}} \binom{k_{(i)}}{h_{(i)}} \Big(-e^{-\mu t i}\Big)^{h_{(i)}} \nonumber\\
	& = e^{-\lambda t} \sum_{n\ge0} \frac{1}{n!} \sum_{\substack{k_1,\dots,k_I=0\\k_1+\dots+k_I=n}0}^n \binom{n}{k_1,\dots, k_I} \prod_{i\in\mathcal{I}} \bigg(\frac{\lambda_i}{i}\bigg)^{k_{(i)}}  \sum_{h_1 = 0}^{k_1} \binom{k_1}{h_1} (-1)^{h_1}\cdots \sum_{h_I = 0}^{k_I} \binom{k_I}{h_I} (-1)^{h_I} \nonumber\\
	&\ \ \ \times \frac{e^{-\mu t\sum_{j\in\mathcal{I}} jh_{(j)} } }{\mu^n}. \label{trasformataLaplaceIntegraleSkellamGeneralizzato}
\end{align}

Now, to derive the probability density \eqref{probabilitaEsplicitaIntegraleSkemllamGeneralizzato}, in light of the linearity of the (bilateral) Laplace transform, it sufficies to invert the quantity in the last line of \eqref{trasformataLaplaceIntegraleSkellamGeneralizzato}, i.e. a function of the form $e^{-\mu H}/\mu^n, \ n\ge0$. It is easy to prove that its inverse is $g(x) = \mathds{1}_{(H,\infty)}(x) (x-H)^{n-1}/(n-1)!, \ x\in\mathbb{R}$. Indeed, with $\mu\in\mathbb{R},$
\begin{equation}\label{trasformazioneLaplaceFunzionePerLeggeEsplicitaSkellam}
	\int_{-\infty}^\infty e^{-\mu x} \mathds{1}_{(H,\infty)}(x) \frac{(x-H)^{n-1}}{(n-1)!}\dif x = e^{-\mu H}	\int_{-\infty}^0 e^{-\mu y} \frac{y^{n-1}}{(n-1)!}\dif x = \frac{e^{-\mu H}}{\mu^n},
\end{equation}
which completes the proof.
\end{proof}

In the case of $\mathcal{I} = \{1\}$ some calculation yields the distribution of the integral of the homogeneous Poisson process, see for instance Section 3 of \cite{SVw2007}. %Note that formula \eqref{leggeIntegralePoissonOmogeneo} can also be derived by means of an induction argument, see Appendix \cite{dimostrazioneLeggeIntegralePoissonOmogeneoRicorrente}.

\begin{remark}[Skellam process of order $K$]
In the case of $M$ being a Skellam process of order $K>0$ (i.e. $\mathcal{I} = \{-K,\dots,-1,1,\dots,K\}$), formula \eqref{probabilitaEsplicitaIntegraleSkemllamGeneralizzato} simplifies as follows, $x\in\mathbb{R}$,
\begin{align}
	P\bigg\{&\int_0^t M(s)\dif s\in \dif x\bigg\} \nonumber\\
	&=  e^{-\lambda t} \sum_{n\ge0} \frac{t^n}{n!} \sum_{\substack{k_{-K},\dots,k_K=0\\k_{-K}+\dots+k_K = n}}^n \binom{n}{k_{-K},\dots, k_{K}} \prod_{i=1}^K \Bigg( \frac{\lambda_{-i}^{k_{-i}} \lambda_{i}^{k_i}}{i^{k_{-i}+k_i}} \sum_{h_i = 0}^{k_{-1}+k_i} \binom{k_{-i}+k_i}{h_i} (-1)^{(k_{-i}+k_i-h_i)} \Bigg)\nonumber \\
	&\ \ \ \times\mathds{1}_{(t\sum_{j=1}^K j(k_j-h_j) ,\infty)}(x) \frac{\big(x-t\sum_{j=1}^K j(k_j-h_j)\big)^{n-1}}{(n-1)!}\dif x, \label{probabilitaEsplicitaIntegraleSkemllamOrdineK}
\end{align}
see Appendix \ref{appendiceLeggeIntegraleSkellamOrdineK} for the details. Note that the case with $K=1$ appeared in \cite{X2018}.
\end{remark}

%we define the compound Poisson process $M(t) = \sum_{k=1}^{N(t)} X_k$ (which is compound Poisson since Proposition \ref{proposizioneCompostoDiCompostoProcessoConteggio}). Here we derive an explicit formula for the probability distribution of the integral of $M$.



\subsection{Renewal case}

We begin by proving some general results involving the Fourier and the Laplace transforms. These are then applied to derive some results about the integrals of compound renewal processes in the Fourier-Laplace domain.
For suitable $f$ and $\mu>0$, hereafter we use the following notation, $\mathcal{L} \{f\}(\mu) =\mathcal{L} f(\mu)= \int_0^\infty e^{-\mu t} f(t)\dif t$.

\begin{lemma}\label{lemmaTrasformataFourierLaplaceSommeLegateRinnovo}
	Let $N$ be a renewal process with waiting times $X\sim F$ and $\{A_n\}_{n\ge1}$ be a sequence of real random variables independent of $N$ (and its waiting times),
	\begin{itemize}
		\item[($i$)] with $\gamma\in\mathbb{R}$,
		\begin{equation}\label{trasformataCompoundRinnovo}
			\mathbb{E} \exp\Biggl( i\gamma \sum_{k=1}^{N(t)} X_k A_k \Biggr) = \sum_{n\ge0}\mathbb{E} \Biggl[\conv_{k=1}^n \Big( e^{i\gamma A_k\sbullet} F(\sbullet) \Big)\ast 1 (t)- \conv_{k=1}^{n} \Big( e^{i\gamma A_k\sbullet}F(\sbullet) \Big)\ast F \ast 1 (t) \Biggr],
		\end{equation}
		where $hF\ast g (t)= \int_0^t g(t-x)h(x)F(\dif x)$ for suitable functions $h,g$ and distribution $F$.
		\item[($ii$)] If $X$ is absolutely continuous with density $f$ and $\mu>0$,
		\begin{equation}\label{trasformataLaplaceFourierCompoundRinnovo}
			\mathcal{L} \left\{ \mathbb{E} \exp\Biggl( i\gamma \sum_{k=1}^{N(t)} X_k A_k \Biggr) \right\} (\mu) = \frac{1-\mathcal{L}f(\mu)}{\mu}\sum_{n\ge0}\mathbb{E} \prod_{k=1}^n \mathcal{L}f(\mu+i\gamma A_k).
		\end{equation}
	\end{itemize}
\end{lemma}

\begin{proof}
($i$) Let $\mathcal{A}_n = \text{Supp}(A_1)\times \dots \times\text{Supp}(A_n)$, then, 
\begin{align*}
	\mathbb{E} &\exp\Biggl( i\gamma \sum_{k=1}^{N(t)} X_k A_k \Biggr) \\
	&= \sum_{n\ge0} P\{N(t) = n \}\int_\mathcal{A} P\{A_1\in \dif a_1,\dots, A_n\in\dif a_n\}\int_0^t e^{i\gamma a_1 x_1} F(\dif x_1)\int_0^{t-x_1} e^{i\gamma a_2x_2}F(\dif x_2)\cdots \\
	& \ \ \ \times \cdots \int_0^{t-x_1-\dots - x_{n-1}} e^{i\gamma a_n x_n}P\{X_{n+1}>t-x_1-\dots-x_n\}F(\dif x_n)\\
	& = \sum_{n\ge0} \int_\mathcal{A} P\{A_1\in \dif a_1,\dots, A_n\in\dif a_n\}\int_0^t e^{i\gamma a_1x_1} F(\dif x_1)\int_0^{t-x_1} e^{i\gamma a_2x_2}F(\dif x_2)\cdots \\
	& \ \ \ \times \cdots \int_0^{t-x_1-\dots - x_{n-1}}e^{i\gamma a_n x_n}F(\dif x_n)\left(1-\int_0^{t-x_1-\dots-x_n}F(\dif x_{n+1})\right)
\end{align*}
which coincides with \eqref{trasformataCompoundRinnovo}.

($ii$) Assuming that $X$ has density $f$,
\begin{align*}
		\mathcal{L} \Bigg\{& \mathbb{E} \exp\Biggl( i\gamma \sum_{k=1}^{N(t)} X_k A_k \Biggr) \Bigg\} (\mu) \\
		&= \sum_{n\ge0}\mathbb{E} \Biggl[ \mathcal{L}\left\{ \conv_{k=1}^n \Big( e^{i\gamma A_k\sbullet}F(\sbullet) \Big)\ast 1 \right\}(\mu) - \mathcal{L}\left\{ \conv_{k=1}^n \Big( e^{i\gamma A_k\sbullet} F(\sbullet)\Big)\ast F\ast 1 \right\}(\mu)\Biggr]\\
		& =  \sum_{n\ge0}\mathbb{E} \left[ \frac{1}{\mu} \prod_{k=1}^n \mathcal{L}f(\mu+i\gamma A_k) - \frac{\mathcal{L}f(\mu)}{\mu} \prod_{k=1}^n \mathcal{L}f(\mu+i\gamma A_k) \right],
\end{align*}
which coincides with \eqref{trasformataLaplaceFourierCompoundRinnovo}.
\end{proof}

%Lemma \ref{lemmaTrasformataFourierLaplaceSommeLegateRinnovo} yields the following statement.
Now we obtain the following results for the integral of a compound renewal process.

\begin{theorem}\label{teoremaTrasformateIntegraleCompoundRinnovo}
	Let $M\sim CR(X,Y)$. 
	\begin{itemize}
		\item[($i$)] With $\gamma \in\mathbb{R}$ and $H_Y(\gamma) = \mathbb{E}e^{i\gamma Y}$,
		\begin{align}
			\mathbb{E}& \exp\left( i\gamma \int_0^t M(s)  \dif s\right ) \nonumber\\
			&= \sum_{n\ge0} \mathbb{E}  \Biggl[\conv_{k=1}^n \Big(  e^{i\gamma \sum_{h=k}^nY_h\sbullet}F(\sbullet)\Big)\ast 1 (t)- \conv_{k=1}^{n} \Big( e^{i\gamma \sum_{h=k}^nY_h\sbullet}F(\sbullet) \Big)\ast F \ast 1 (t) \Biggr] \label{fourierIntegraleCompoundRinnovoOsservandoX} \\
			& = \sum_{n\ge0} \mathbb{E} \Bigg[ \mathds{1}(X> t-S_n ) \prod_{k=1}^n H_Y\big(\gamma (t-S_k)\big)\Bigg] . \label{fourierIntegraleCompoundRinnovoOsservandoY}
		\end{align}
		\item[($ii$)] If $X$ is absolutely continuous with density $f$ and $\mu>0$,
		\begin{align}
		\int_0^t e^{-\mu t} \mathbb{E} \exp\Bigg( i\gamma\bigg(t\sum_{k=1}^{N(t)} Y_k &- \int_0^t M(s)  \dif s\bigg) \Bigg) \dif t = \frac{1-\mathcal{L}f(\mu)}{\mu}\sum_{n\ge0}\mathbb{E} \prod_{k=1}^n \mathcal{L}f\bigg(\mu+i\gamma \sum_{h=k}^n Y_k\bigg). \label{fourierLaplaceIntegraleCompoundRinnovo}
	\end{align}
	\end{itemize}
\end{theorem}

\begin{proof}
	Points ($i$) and ($ii$) of Lemma \ref{lemmaTrasformataFourierLaplaceSommeLegateRinnovo} respectively yield forumla \eqref{fourierIntegraleCompoundRinnovoOsservandoX} and formula \eqref{fourierLaplaceIntegraleCompoundRinnovo}.
	Formula \eqref{fourierIntegraleCompoundRinnovoOsservandoY} follows by means of \eqref{funzioneCaratteristicaIntegraleGenerale} and by considering that in the renewal case $P\{X_1\in \dif x_1, \dots, X_n\in \dif x_n, N(t)=n\} = P(X>t-x_1-\dots-x_n) \prod_{k=1}^n P(X_k\in\dif x_k),\ n\ge1$.
\end{proof}


We point out that Theorem \ref{teoremaTrasformateIntegraleCompoundRinnovo} generalizes Theorem 4.1 of \cite{SVw2007}. Indeed, if $Y=1$ a.s. formula \eqref{fourierLaplaceIntegraleCompoundRinnovo} concerns the integral of the renewal process $N$ and it coincides with equation (7) of \cite{SVw2007}.




%\section{Maximum likelihood inference}


%%%%%%%%%%%%%%%%%%%%%%%%%%

%\subsection*{\large{Declarations}}
%
%\textbf{Ethical Approval.} This declaration is not applicable.
% \\
%\textbf{Competing interests.}  The authors have no competing interests to declare.
%\\
%\textbf{Authors' contributions.} Both authors equally contributed in the preparation and the writing of the paper.
% \\
%\textbf{Funding.} The authors received no funding.
% \\
%\textbf{Availability of data and materials.} This declaration is not applicable.

%\subsection*{Acknowledgments}

% ---------------------------------------------------------------------------------------------------------

%% The Appendices part is started with the command \appendix;
%% appendix sections are then done as normal sections

 \appendix

 
 \section{Proof of formulas \eqref{formulazioneGeneraleIntegraleProcessoTempoContinuo} and \eqref{formulazioneGeneraleIntegraleProcessoCompound}}\label{appendiceCalcoloEsplicitoIntegrali}
 Formula \eqref{formulazioneGeneraleIntegraleProcessoTempoContinuo}.
\begin{align*}
	\int_0^t X_{N(s)} f(s)\dif s &=  \sum_{n=0}^\infty  X_n \int_{0}^t \mathds{1}\big(N(s) =n\big) f(s)\dif s \\
	&= \sum_{n=0}^\infty  X_n \int_{0}^t \mathds{1}\big(S_n \le s < S_{n+1}\big) f(s)\dif s \\
	&= \sum_{n=0}^\infty  X_n \int_{\min\{S_n, t\}}^{\min\{t, S_{n+1}\}} f(s)\dif s \\
	&= \sum_{n=0}^{N(t)-1}  X_n \int_{S_n}^{S_{n+1}} f(s)\dif s + X_{N(t)} \int_{S_{N(t)}}^t f(s)\dif s,
\end{align*}
where in the first step we exchanged series and integral in light of Condition II.

Formula \eqref{formulazioneGeneraleIntegraleProcessoCompound}.
\begin{align*}
	\int_0^t  \sum_{k=1}^{N(s)} Y_k f(s)\dif s &=  \sum_{n=0}^{N(t)-1}  \sum_{k=1}^{n} Y_k \int_{S_n}^{S_{n+1}} f(s)\dif s + \sum_{k=1}^{N(t)} Y_k \int_{S_{N(t)}}^t f(s)\dif s \\
	& = \sum_{k=1}^{N(t)-1} Y_k  \int_{S_k}^{S_{N(t)}} f(s)\dif s  + \sum_{k=1}^{N(t)} Y_k \int_{S_{N(t)}}^t f(s)\dif s\\
	& = \sum_{k=1}^{N(t)} Y_k   \int_{S_k}^t f(s)\dif s
\end{align*}
where in the last equation we used that in the first sum we can add the term $k=N(t)$ since it is equal to $0$.

%\section{Induction proof of Formula \eqref{leggeIntegralePoissonOmogeneoRicorrente}}\label{dimostrazioneLeggeIntegralePoissonOmogeneoRicorrente}

\section{Proof of Lemma \ref{lemmaIncrementiIntegraleProcesso}}\label{appendiceDimostrazioneIncrementiIntegrale}
($i$) For $\gamma\in\mathbb{R}$ and $0\le t_1<t_2\le T$,
\begin{align}
	\mathbb{E}e^{i\gamma \big(I(t_2)-I(t_1)\big)}  &=  \mathbb{E} \exp\Bigg( i\gamma \int_{t_1}^{t_2} \Big(X(s)\pm X(t_1) \Big) \dif s \Bigg) \nonumber\\
	& = \mathbb{E} \exp\Bigg( i\gamma \int_{t_1}^{t_2} \Big(X(s)- X(t_1) \Big) \dif s  +  i\gamma (t_2 - t_1) X(t_1)\Bigg) \nonumber
\end{align}
and the independence of the increments yields \eqref{incrementoIntegraleProcessoStocasticoIncrementiIndipendenti}.
\\
Point ($ii$) readily follows from ($i$) by observing that $\{X(s+t_1)-X(t_1)\}_{s\ge0} \stackrel{d}{=}X$ on $\mathcal{D}[0,T]$ because of ($ii$) of Lemma \ref{lemmaUguaglianzaProcessiSpazioFunzioniCadlag}. Hence, the integral of the increment process is equal in distribution to the integral of $X$.

\section{Proof of Formula \eqref{probabilitaEsplicitaIntegraleSkemllamOrdineK}}\label{appendiceLeggeIntegraleSkellamOrdineK}
We recover the steps of the calculation of the bilateral Laplace transform in \eqref{trasformataLaplaceIntegraleSkellamGeneralizzato} which can be written as follows, with $\mu\in\mathbb{R}$,
\begin{align}
	\mathbb{E}&\exp\bigg( -\mu \int_0^t M(s)\dif s\bigg) \nonumber\\
	%& = e^{-\lambda t}\sum_{n\ge0} \frac{1}{n!} \bigg( \sum_{i\in\mathcal{I}} \lambda_i t \text{sgn}(i) \frac{1-e^{-\mu t i}}{\mu}\bigg)^n\nonumber\\
	& = e^{-\lambda t} \sum_{n\ge0} \frac{\mu^{-n}}{n!}  \Bigg( \sum_{\substack{i=-K\\ i\not=0}}^K \frac{\lambda_i}{i} \Big(1-e^{-\mu t i}\Big)\Bigg)^n\nonumber\\
	& = e^{-\lambda t} \sum_{n\ge0} \frac{1}{n!} \bigg(\frac{t}{\mu}\bigg)^n \Bigg(\sum_{i=-1}^{-K} \frac{\lambda_i}{|i|} (e^{-\mu t i}-1)  +\sum_{i=1}^K \frac{\lambda_i}{i} \frac{e^{\mu t i}-1}{e^{\mu t i}}\Bigg)^n \nonumber\\
	& =  e^{-\lambda t} \sum_{n\ge0} \frac{1}{n!} \bigg(\frac{t}{\mu}\bigg)^n \sum_{\substack{k_{-K},\dots,k_K=0\\k_{-K}+\dots+k_K = n}}^n  \binom{n}{k_{-K},\dots,k_K} \prod_{i=1}^K \frac{\lambda_{-i}^{k_{-i}} \lambda_i^{k_i}}{i^{k_{-i}+k_i}} \frac{(e^{\mu t i}-1)^{k_{-i}+k_i}}{e^{\mu t i k_i}}\nonumber\\ 
	& = e^{-\lambda t} \sum_{n\ge0} \frac{1}{n!} \bigg(\frac{t}{\mu}\bigg)^n \sum_{\substack{k_{-K},\dots,k_K=0\\k_{-K}+\dots+k_K = n}}^n \Bigg(\prod_{i=1}^K \frac{\lambda_{-i}^{k_{-i}}\lambda_{i}^{k_i}}{i^{k_{-i}+k_i}}\Bigg) \sum_{h_1 = 0}^{k_{-1}+k_1}\binom{k_{-1}+k_1}{h_1}\cdots \sum_{h_K=0}^{k_{-K} + k_K} \binom{k_{-K} + k_K}{h_K} \nonumber\\
	& \ \ \ (-1)^{\sum_{j=1}^K(k_{-j}+k_j-h_j)}\times e^{-\mu t \sum_{j=1}^K j (k_j - h_j)}. \label{trasformataLaplaceIntegraleSkellamOrdineK}
\end{align}
By means of \eqref{trasformazioneLaplaceFunzionePerLeggeEsplicitaSkellam} we obtain that the inverse of \eqref{trasformataLaplaceIntegraleSkellamOrdineK} is \eqref{probabilitaEsplicitaIntegraleSkemllamOrdineK}.
% \section{}
% \label{}

%% For citations use: 
%%       \cite{<label>} ==> Jones et al. [21]
%%       \citep{<label>} ==> [21]
%%

%% If you have bibdatabase file and want bibtex to generate the
%% bibitems, please use
%%
%%  \bibliographystyle{elsarticle-num-names} 
%%  \bibliography{<your bibdatabase>}

%% else use the following coding to input the bibitems directly in the
%% TeX file.
\footnotesize{
\begin{thebibliography}{00}
%% \bibitem[Author(year)]{label}
%% Text of bibliographic item
%
%\bibitem{A1838}
%Airy, G.B. (1838), On the intensity of light in the neighbourhood of a caustic, Trans. Camb. Phil. Soc. 6, 379--402.
%
%\bibitem{A2017}
%Ansari, A. (2017) Riesz fractional derivatives of solutions of differentual equation $x^{(4)} + xy = 0$, Studia Scientiarum Mathematicarum Hungarica 54(2), 205--220. 
%
%\bibitem{AA2014}
%Ansari, A., Askari, H. (2014), On fractional calculus of $A_{2n+1}(x)$ function, Applied Mathematics and Computation 232, 487--497.
%
%\bibitem{A2009}
%Applebaum, D., L\'{e}vy Processes and Stochastic Calculus, 2nd Edition, Cambridge University Press, 2009.
%\bibitem{AA2020}
%Askari, H., Ansari, A.  (2020), On Mellin transforms of solutions of differential equation $\chi^{(n)}(x)+\gamma_nx\chi(x)=0$, Analysis and Mathematical Physics 10, no.4, 57.
%
%\bibitem{AA2024}
%Askari, H., Ansari, A., Stokes phenomenon and representations of M-wright functions in uniform asymptotic expansions with many coalescing saddle points: order $\frac{1}{n}$, Preprint (2024).
%
%%\bibitem{B2008}
%%Beghin, L. (2008), Pseudoprocesses governed by higher-order fractional differential equations, Electronic Journal of Probability 13, 467--485.
%

\bibitem{BCS2026}
Beghin, L., Cristofaro, L., Scalas, E. (2026), Shot-noise processes with logarithmic response function and their scaling limits, arXiv:2602.03503.

\bibitem{BM2012}
Beghin, L., Macci, C. (2012), Alternative forms of compound fractional Poisson processes, Abstr. Appl. Anal. 747503.

\bibitem{BM2014}
Beghin, L., Macci, C. (2014), Fractional discrete processes: compound and mixed Poisson representations, J. Appl. Probab. 51(1), 19--36.
%
%\bibitem{BKOS2007}
%L. Beghin, L., Kozachenko, Y., Orsingher, E., Sakhno, L. (2007), On the solutions of linear odd-order heat-type equations with random initial conditions, J. Stat. Phys. 127(4), 721--739.
%
%\bibitem{B1999}
%Billingsley, P., Convergence in probability measure, Second Edition, Wiley Series in Statistics and Probability, John Wiley and Sons, Inc., 1999.
%\bibitem{BM2015}
%Bonaccorsi, S., Mazzucchi, S. (2015), High order heat-type equations and random walks on the complex plane, Stoch. Process. Appl. 125(2), 797--818.
%
\bibitem{B2022}
Borovkov, A.A. (2022), Compound Renewal Process, Cambridge University Press.
%
%%\bibitem{B1999}
%%Brooks, E.A. (1999), Probabilistic methods for linear reaction-hyperbolic system with constant coefficients, The Annals of Applied Probability 9, 719--731.
%
\bibitem{BS2024}
Buchak, K., Sakhno, L. (2024), Generalized fractional calculus and some models of generalized counting processes, Modern Stoch. Theory Appl. 11(4), 439--458.
%\bibitem{C2017}
%Chen, Z.Q. (2017), Time fractional equations and probabilistic representation, Chaos, Solitons Fractal 102, 168--174.
%\bibitem{C2022}
%Cinque, F. (2022), On the sum of independent generalized Mittag-Leffler random variables and the related fractional processes, Stochastic Analysis and Applications 40(1), 103--117.
%
%%\bibitem{CO2020}
%%Cinque, F., Orsingher, E. (2020), On the distribution of the maximum of the telegraph process, Theory of Probability and Mathematical Statistics 102, 73--95.
%% 
%%\bibitem{C2022}
%%Cinque, F. (2022), A note on the conditional probabilities of the telegraph process, Statistics and Probability Letters 185, 109431.
%
%%\bibitem{CO2021a}
%%Cinque, F., Orsingher, E. (2021), On the exact distribution of the maximum of the asymmetric telegraph process, Stochastic Processes and their Applications 142, 601--633.
%
%%\bibitem{CO2021b}
%%Cinque, F., Orsingher, E., Stochastic dynamics of generalized planar random motions with orthogonal directions, Preprint (2021) arXiv:2108.10027.
%
%\bibitem{CO2024}
%Cinque, F., Orsingher, E. (2024), Analysis of Fractional Cauchy problems with some probabilistic applications, J. Math. Anal. Appl. 536, 128188.

\bibitem{CO2025}
Cinque, F., Orsingher, E. (2025), Higher-order fractional equations and related time-changed pseudo-processes, J. Math. Anal. Appl. 543, 129026.

\bibitem{CO2026}
Cinque, F., Orsingher, E. (2026), Interacting point processes, Journal of Theoretical Probability 39:47.

\bibitem{CO202}
Cinque, F., Orsingher, E., Point processes of the Poisson-Skellam family, Preprint (2026), arXiv:2504.07672.
%\bibitem{DF1965}
%Daletsky, Y.L., Fomin, S.V. (1965), Generalized measures in function spaces, Theory Probab. Appl. 10(2), 304--316.
%
%\bibitem{D2006}
%Debbi, L. (2006), Explicit solutions of some fractional partial differential equations via stable subordinators, J. Appl. Math. Dtoch. Anal. 5:093502.
%
%%\bibitem{Dg2014}
%%De Gregorio, A., (2014), A family of random walks with generalized Dirichlet steps, J. Math. Phys. 55, 023302.
%
%%\bibitem{DgI2021}
%%De Gregorio, A., Iafrate, F. (2021), Telegraph random evolutions on a circle, Stochastic Processes and their Applications 141, 79--108.
%
%\bibitem{DK2024}
%Dhillon, M., Kataria, K.K. (2024), On the superposition and thinning of generalized counting processes, Stochastic Analysis and Applications 42(6), 1110--1136
%.
%%\bibitem{Dc2002}
%%Di Crescenzo, A. (2002), Exact transient analysis of a planar motion with three directions, Stochastics and Stochastics Reports 72(3-4), 175--189.

\bibitem{DcMZ2015}
Di Crescenzo, A., Martinucci, B., Zacks, S. (2015), Compound Poisson process with a Poisson subordinator, Journal of Applied Probability 52(2), 360--374.

\bibitem{DcMM2016}
Di Crescenzo, A., Martinucci, B., Meoli, A. (2016), A fractional counting process and its connection with the Poisson process, ALEA, Lat. Am. J. Probab. Math. Stat. 13, 291--307.




%%\bibitem{EG2015}
%%Elgeti, J., Gompper, G. (2015), Run-and-Tumble dynamics of self-propelled particles in confinement, Eur. Phy. Lett. 109(5), 58003.
%%
%%\bibitem{FVk1992}
%%Foong, S.K., Van Kolck, U. (1992), Poisson random walk for solving wave equations, Progress of Theoretical Physics 87, 285--292.
%%
\bibitem{GMn1971}
Gani, J., McNeil, D.R. (1971), Joint distributions of random variables and their integrals for certain birth-death and diffusion processes, Adv. Appl. Prob. 3, 339--352.
%\bibitem{GOP2015}
%Garra, R., Orsingher, E., Polito, F. (2015), State-dependent fractional point processes, J. Appl. Probab. 52(1), 18--36.
%%\bibitem{GO2016}
%%Garra, R., Orsingher, E. (2016), Random flights related to the Euler-Poisson-Darboux, Markov Processes and Related Fields 22, 87--110.
%
%\bibitem{GOS2017}
%Garra, R., Orsingher, E., Scavino, M. (2017), Some probabilistic properties of fractional point processes, Stochastic Analysis and Application 35(4), 701--718.
%%\bibitem{G1951}
%%Goldstein, S. (1951), On diffiusion by discontinuous movements and the telegraph equation, Quart. J. Mech. Appl. Math. 4, 129--156.
%
%\bibitem{GHPD2013}
%G\'{o}rska, K., Horzela, A., Penson, K.A., Dattoli, G. (2013), The higher-order heat-type equations via signed L\'{e}vy stable and generalized Airy functions, Journal of Physics A: Mathematical and theoretical 46, 425001.
%
%%\bibitem{GR1980}
%%Gradshteyn, I.S. and Ryzhik, I.M. (1980), Tables of Integrals, Series and Products. Academic Press, New York.
%
%%\bibitem{G1953}
%%Grosjean, C.C. (1953), Solution of the non-isotropic random flight problem in the $k$-dimensional space, Physica 19, 1--12, 29--45.
%
%
%\bibitem{GKL2020}
%Gupta, N., Kumar, A., Leonenko, N. (2020), Skellam type processes of order $k$ and beyond, Entropy 22, 1193.

%\bibitem{G2005}
%Gut, A., Probability: A Graduate Course, Springer Texts in Statistics, Springer, 2005.
%%\bibitem{HMSS2020}
%%Hartmann, A.K., Majumdar, S.N., Schawe, H., Schehr, G. (2020), The convex hull of the run-and-tumble particle in a plane, J. Stat. Mech. 053401.
%
%\bibitem{HJK2007}
%Hwang, Y., Kim, J., Kweon, I. (2007), Sensor noise modeling using the Skellam distribution: Application to the color edge detection, Proceedings of the IEEE Conference on Computer Vision and Pattern Recognition, Minneapolis, MN, USA, 17–22 June 2007; pp. 1–8.
%\bibitem{H1978}
%Hochberg, K.J. (1978), A signed measure on path space related to wiener measure, Ann. Probab. 6(3), 433--458.
%
%%\bibitem{I2001}
%%Iacus, S.M. (2001), Parametric estimation for the standard and geometric telegraph process observed at discrete times, Stat. Infer. Stoch. Process. 11, 249--263
%%
%%%\bibitem{}
%%%Kac, M. (1974), A stochastic model related to the telegrapher's equation, \textit{Rocky Mountain Journal of Math.}, \textbf{4}, 497--509.
%%
%\bibitem{KN2008}
%Karlis, D.; Ntzoufras, I. (2008), Bayesian modeling of football outcomes: Using the Skellam’s distribution for the goal difference, IMA J. Manag. Math. 20, 133–145.
%
\bibitem{KM2021}
Kataria, K.K., Khandakar, M. (2021), Convoluted fractional Poisson process, ALEA, Lat. Am. J. Probab. Math. Stat. 18, 1241--1265.

\bibitem{KK2023}
Kataria, K.K., Khandakar, M. (2023), Some compound fractional Poisson processes. Fractal Fract. 7, 15.
%
\bibitem{KK2024}
Kataria, K.K., Khandakar, M. (2024), Fractional Skellam process of order $k$, J. Theor. Probab. 37, 1333--1356.

\bibitem{KV2019}
Kataria, K.K., Vellaisamy, P. (2019), On distributions of certain state dependent fractional point processes, J. Theor. Probab. 32(3), 1554--1580.
%
%%\bibitem{K2018}
%%Kolesnik, A.D. (2018), Slow diffusion by Markovian random flights, Theory of Stochastic Processes 15(1), 33--39.
%%
%%\bibitem{KO2001}
%%Kolesnik, A.D., Orsingher, E. (2001), Analysis of finite velocity planar random motion with reflection, Theory of Probability and its Applications 46(1), 132--140.
%%
%%\bibitem{KO2005}
%%Kolesnik, A.D., Orsingher, E. (2005), A planar random motion with an infinite number of directions controlled by the damped wave equation, Journal of Applied Probability 42(4), 1168--1182.
%
%%\bibitem{KT1998}
%%Kolesnik, A.D., Turbin, A.F. (1998), The equation of symmetric Markovian random evolution in a plane, Stochastic Processes and their Applications 75, 67--87. 
%

%\bibitem{K1960}
%Krylov, V.Yu. (1960), Some properties of the distribution corresponding to equation $\frac{\partial u}{\partial t} = (-1)^{q+1}\frac{\partial^{2q} u}{\partial x^{2q}}$, Dokl. Akad. Nauk SSSR 132(6), 1254--1257.

\bibitem{K2011}
Kochubei, A.N. (2011), General fractional calculus, evolution equations, and renewal processes, Integral Equ. Oper. Theory 71(4), 583--600.

%\bibitem{L2003}
%Lachal, A. (2003), Distributions of sojourn time, maximum and minimum for pseudo-processes governed by higher-order heat-type equations, Elect. J. Probab. 8:1--53.
%
%%\bibitem{L2006}
%%Lachal, A. (2006), Cyclic random motions in $\mathbb {R}^d$-space with $n$ directions, ESAIM: Probability and Statistics 10, 277--316.
%
%\bibitem{L2007}
%Lachal, A. (2007), First hitting time and place, monopoles and multipoles for pseudo-processes driven by the equation $\partial u/\partial t = \pm \partial^N u/ \partial x^N$, Electr. J. Probab. 12:300-353.
%
%\bibitem{L2012}
%Lachal, A. (2012), A survey on the pseudo.process driven by the high-order heat-type equation $\partial /\partial t =\pm \partial^N/\partial x^N$ concerning the hitting and sojourn times, Methodol. Comput. Appl. Probab. 14, 549--566.
%
%%\bibitem{LLO2006}
%%Lachal, A., Leorato, S., Orsingher, E. (2006), Minimal cyclic random motion in ${R}^{n}$ and hyper-Bessel functions, Annales de l'I.H.P. Probabilités et statistiques 42 (6), 753--772.
%
\bibitem{L2003}
Laskin, N. (2003), Fractional Poisson process, Commun. Nonlinear Sci. Numer. Simul. 8(3--4), 201--213.

%%\bibitem{LO2004}
%%Leorato, S., Orsingher, E. (2004), Bose-Einstein-type statistics, order statistics and planar random motions with three directions, Adv. Appl. Prob. 36, 937--970.

\bibitem{MV2019}
Maheshwari, A., Vellaisamy, P. (2019), Non-homogeneous space-time fractional Poisson processes, Stochastic Analysis and Applications 37(2), 137--154.

\bibitem{MGS2004}
Mainardi, F., Gorenflo, R., Scalas, E. (2004), A fractinoal generalization of the Poisson processes, Vietnam Journal of Mathematics 32, 52--64.%%\bibitem{LR2014}
%%Lopez, O., Ratanov, N. (2014), On the asymmetric telegraph processes, J. Appl. Prob. 51, 569--589.
%
%\bibitem{M2010}
%Mainardi, F., Fractional Calculus and Waves in Linear Viscoelasticity, Imperial College Press, London, 2010.
%
%\bibitem{GKMR2014}
%Gorenflo, R., Kilbas, A. A., Mainardi, F., Rogosin, S. V., Mittag-Leffler Functions. Related Topics and Applications, Heidelberg: Springer, 2014.
%
%\bibitem{MO2023}
%Marchione, M.M., Orsingher, E. (2023), Stable distributions and pseudo-processes related to fractional Airy functions, Stochastic Analysis and Applications 42(2), 435--450.

\bibitem{MS2025}
Martinuccci, B., Spina, S. (2025), On a finite-velocity random motion governed by a modified Euler-Poisson-Darboux equation, Electron. J. Probab. 30, 1--28.

%%\bibitem{MPW1993}
%%Masoliver, J., Porra, J.M., Weiss, G.H. (1993), Some two and three-dimensional persistent random walks, Physica A: Statistical Mechanics and its Applications 193(3-4), 469-482
%
%%\bibitem{ML2020}
%%Masoliver, J., Lindenberg, K. (2020), Two-dimensional telegraphic processes and their fractional generalizations, physical review E 101(1), 012137.
%
%%\bibitem{M2021}
%%Masoliver, J. (2021), Telegraphic transport processes and their fractional generalization: a review and some extensions, Entropy 23(3), 364.
%
\bibitem{Mn1970}
McNeil, D.R. (1970), Integral functional of birth and death processes and related limiting distributions, Ann. Math. Statist. 41, 480--485.

\bibitem{MNV2011}
Meerschaert, M.M., Nane, E., Vellaisamy, P. (2011), The fractional Poisson process and the inverse
stable subordinator, Electron. J. Probab. 16(59), 1600--1620.
%%\bibitem{MADB2021}
%%Mertens, K., Angelani, L., Di Leonardo, R., Bocquet, L. (2012), Probability distributions for the run-and-tumble bacterial dynamics: An analogy to the Lorentz model, The European Physical Journal 35, 84.
%
%%\bibitem{MLdMS2020}
%%Mori, F., Le Doussal, P., Majumdar, S.N., Schehr, G. (2020), Universal Properties of a Run-and-Tumble Particle in Arbitrary Dimension, Phys. Rev. E 102, 042133.
%
\bibitem{N1966}
Naddor, E. (1966), Inventory Systems, Wiley, New York.
%%\bibitem{O1990}
%%Orsingher, E. (1990), Probability law, flow function, maximum distribution of wave-governed random motions and their connections with Kirchoff's laws, Stochastic Processes and their Applications 34(1), 49--66.
%%\bibitem{O1986}
%%Orsingher, E. (1986), A planar random motion governed by the two-dimensional telegraph equation, Journal of Applied Probability 23(2), 385--397.
%%
%%\bibitem{O2000}
%%Orsingher, E. (2000), Exact joint distribution in a model of planar random motion, Stochastics and Stochastics Reports 69, 1--10.
%%
%%\bibitem{O2002}
%%Orsingher, E. (2002), Bessel functions of third order and the distribution of cyclic planar random motion with three directions, Stochastics and Stochastics Reports 74, 617--631.
%
%\bibitem{ODo2012}
%Orsingher, E., D'Ovidio, M. (2012), Probabilistic representation of fundamental solutions to $\frac{\partial u}{\partial t} = \kappa_m\frac{\partial^m u}{\partial x^m}$, Electronic Communications in Probability 17, 1--12.
%
%%\bibitem{ODg2007}
%%Orsingher, E., De Gregorio, A. (2007), Random flights in higher spaces, Jour. Theor. Probab. 20, 769--806.
%
%%\bibitem{OGZ2020}
%%Orsingher, E., Garra, R., Zeifman, A.I. (2020), Cyclic random motions with orthogonal directions, Markov Processes Relat. Fields 26, 381--402.
%
%%\bibitem{OK1996}
%%Orsingher, E., Kolesnik, A.D. (1996), Exact distribution for a planar random motion model, controlled by a fourth-order hyperbolic equation, Theory of Probability and its Applications 41, 379--386.
%
\bibitem{OP2012}
Orsingher, E., Polito, F. (2012),  The space-fractional Poisson process, Statistics and Probability Letters, 82, 852--858.

\bibitem{OP2013}
Orsingher, E., Polito, F. (2013),  On the integral of the fractional Poisson processes, Statistics and Probability Letters, 83, 1006--1017.
%\bibitem{OT2014}
%Orsingher, E., Toaldo, B. (2014), Pseudoprocesses related to space-fractional higher-order heat-type equations, Stochastic Analysis and Applications 32(4), 619--641.
%
\bibitem{OT2015}
Orsingher, E., Toaldo, B. (2015), Counting processes with Bernstein intertimes and random jumps, J. Appl. Probab. 52, 1028--1044.
%%\bibitem{PDlA2014}
%%Paoluzzi, M., Di Leonardo, R., Angelani, L. (2014), Run-and-tumble particles in speckle fields, J Phys Condens Matter 26(37), 375101.
%
\bibitem{PC2026}
Pedicone, A., Cinque, F. (2026), On the weak convergence of the integral of semi-Markov processes, J. Appl. Probab., 1--29.
%%\bibitem{P1906}
%%Pearson, F.R.S. (1906), A mathematical theory of random migration, Mathematical Contributions to the Theory of Evolution. XV.
%
\bibitem{P1984}
Philippou, A. N. (1984), Poisson and compound Poisson distributions of order $k$ and some of their properties, J. Sov. Math. 27, 3294--3297.
%\bibitem{P1999}
%Podlubny, I., Fractional Differential Equations, Mathematics in Sciences and Engineering , Academic Press, San Diego, 1999.
%
%%\bibitem{P2012}
%%Pogorui, A. (2012), Evolution in multidimensional spaces, Random Oper. Stoch. Equ. 20, 119--126.
%
%%\bibitem{PR2011}
%%Pogorui, A. Rodr\'iguez-Dagnino, R., (2011), Isotropic random motion at finite speed with K-Erlang distributed direction alternations, Journal of Statistical Physics 145, 102--112.
%
%\bibitem{PKS2011}
%Politi, M., Kaizoji, T., Scalas, E. (2011), Full characterization of the fractional Poisson process, EPL. 96(2), 20004.

\bibitem{PS2002}
Pollett, P.K., Stefanov, V.T. (2002), Path integrals for continuous-time random walks, J. Appl. Prob. 39, 901--904.

\bibitem{P1966}
Puri, P.S. (1966), On the homogeneous birth-and-death process and its integral, Biometrika 53, 61--71.
\bibitem{P1968}
Puri, P.S. (1968), Some further results on the birth-death process and its integral, Proc. Camb. Phil. Soc. 64, 141--154.
%\bibitem{R1995}
%Reid, W.H. (1995), Integral representations for products of Airy functions, Zeitschrift f\"{u}r angewandte Mathematik und Physik 46, 159--170.
%\bibitem{R1997}
%Reid, W.H. (1997), Integral representations for products of Airy functions Part 2: Cubic products, Zeitschrift f\"{u}r angewandte Mathematik und Physik 48, 646--655.
%\bibitem{R1997b}
%Reid, W.H. (1997), Integral representations for products of Airy functions, Part 3: Quartic products, Zeitschrift f\"{u}r angewandte Mathematik und Physik 48, 656--664.
%
%%\bibitem{S2001}
%%Samoilenko, I.V. (2001), Markovian evolutions in $\mathbb {R}^n$, Random Oper. Stochastic Equations 9(2), 139--160.
%
%%\bibitem{SBS2020}
%%Santra, I., Basu, U., Sabhapandit, S. (2020), Run-and-Tumble Particles in two dimesions: Marginal position distributions, Phys. Rev. E 101, 062120.

\bibitem{S2013}
Scalas, E. (2013), A class of CTRWs: compound fractional Poisson processes, Fractional Dynamics, World Scientific Publishing: Hackensack, NJ, USA, 353--374.

\bibitem{SSV2012}
Schilling, R. L., Song, R., Vondracek, Z., Bernstein Functions: Theory and Applications, De Gruyter, Berlin, 2012.

%\bibitem{S1950}
%Scorer, R.S. (1950), Numerical evaluation of integrals of the form $I = \int_{x_1}^{x_2} f(x) e^{i\phi(x)} dx$ and the tabulation of the function $Gi(z) = (1/\pi)\int_0^\infty\sin(uz+\frac{1}{3}u^3)du$, Quart. Journ. Mech. and Applied Math., Vol. III, Pt. 1.
%
%\bibitem{SMU2020}
%Sengar, A.S., Maheshwari, A., Upadhye, N.S. (2020), Time-changed Poisson processes of order k. Stoch. Anal. Appl. 38(1), 124--148.
%%\bibitem{S2020}
%%Sevilla, F.J., (2020), Two-dimensional active motions, Phys. Rev. E 101, 022608.
%
%%\bibitem{S2022}
%%Shitikova, M.V. (2022), Fractional operators viscoelastic models in dynamic problems of mechanics of solids: a review, Mechanics of Solids 57, 1--33.
%
\bibitem{S1946}
Skellam, J.G. (1946), The frequency distribution of the difference between two Poisson variates belonging to different populations, J. R. Stat. Soc. (N.S.) 109, 296.
%%\bibitem{S1987}
%%Stadje, W. (1987), The exact probability distribution of a two-dimensional random walk, J. Statist. Phys. 56, 415--435.
%
\bibitem{SVw2007}
Suyono, van der Weide, J.A.M. (2007), Integrated renewal process, J. Indones. Math. Soc. 13(2), 149--159.

\bibitem{T2015}
Toaldo, B. (2015), Convolution-type derivatives, hitting-times of subordinators and time-changed $C_0$-semigroups, Potential Anal. 42, 115--140.
%\bibitem{VS2010}
%Vallée, O., Soares, M., Airy functions and applications to physics, 2nd edition, Imperial College Press, London, 2010.
%
%\bibitem{V2010}
%Varlamov, V. (2010), Integrals involving products of Airy functions, their derivatives and Bessel functions, J. Math. Anal. Appl. 370, 687--702.
%
\bibitem{VK2024}
Vishwakarma, P., Kataria, K.K., (2024), On integrals of birth–death processes at random time, Statistics and Porbability Letters 214, 110204.

\bibitem{W1971}
Watanabe, T. (1971), An integro-differential equation for a compound Poisson process with drift and the integral equation of H.Cramer, Osaka J. Math. 8, 377--383.

\bibitem{X2018}
Xia, W. (2018), On the distribution of running average of Skellam process, Int. J. Pure Appl. Math. 119, 461--473.
%\bibitem{Z1986}
%Zolotarev, V.M. (1986), One-Dimensional Stable Distributions, Translations of Mathematical Monographs, vol.65, American Mathematical Society.

\end{thebibliography}
}

\end{document}

% - - - - -  - - - - -  - - - - -  - - - - -  - - - - -  - - - - -  - - - - -  - - - - - 
% - - - - -  - - - - - FINE DOCUMENTO - - - - -  - - - - - 
% - - - - -  - - - - -  - - - - -  - - - - -  - - - - -  - - - - -  - - - - -  - - - - - 




